\documentclass[../main.tex]{subfiles}
\begin{document}

\chapter{Thread Design Considerations}
\lessoninfo{
  video={P2L4},
  textbook={OSTEP \cite{ostep} ch.~6, 26--28; Silberschatz et al.\ \cite{silberschatz2018} ch.~4; Kerrisk \cite{kerrisk2010} ch.~20--22, 28},
  keywords={hard process state, light process state, lightweight process, red zone, adaptive mutex, reaper thread, interrupt handler, signal handler, signal mask, top half, bottom half, real-time signal, task, NPTL},
}

Lesson~3 introduced kernel-level and user-level threads and the models that
map one kind onto the other, and Lesson~4 the POSIX threads interface
through which applications use them.  This lesson looks inside the
implementation: which data structures the kernel and a user-level thread
library keep, how the two levels coordinate their decisions, and how
interrupts and signals interact with multithreading.  The presentation
follows two papers on the multithreaded kernel of SunOS~5.0
\cite{eykholt1992} and its user-level thread library \cite{stein1992}, and
closes with the way Linux represents threads.

\section{Thread data structures}
\label{sec:l05-structures}

Supporting threads requires data structures at two levels: in the kernel
for kernel-level threads, and in the user-level library for user-level
threads.\source{P2L4, thread data structures: single CPU, at scale; hard
and light process state; rationale for data structures; Eykholt et al.\
\cite{eykholt1992}; OSTEP ch.~26.}  Consider first a process whose
user-level threads all run on a single kernel-level thread (the many-to-one
model of Lesson~3).  The kernel describes the process by its PCB, which
holds the virtual address mappings together with the execution context of
the kernel-level thread: its registers, stack pointer and stack.  The
user-level library keeps a data structure for each user-level thread, with
the thread's identifier, its saved registers and its stack.  By loading one
of these contexts onto the kernel-level thread, the library decides which
user-level thread runs.\margin{The PCB is Lesson~2; here its contents are
split, so that the threads of one process share the address-space part.}

When a process has several kernel-level threads, one PCB can no longer
hold everything, because every kernel-level thread needs its own execution
context.  Replicating the complete PCB for each of them would also
replicate the virtual address mappings, which are the same for all threads
of the process.  The kernel therefore keeps a separate kernel-level thread
data structure with the registers, stack pointer and stack of each
kernel-level thread, and the PCB keeps the information that is common to
the process.\margin{On x86-64 Linux every task carries its own
\SI{16}{\kibi\byte} kernel stack (four pages), while the address space it
points to is shared.  Kernel documentation, \texttt{x86/kernel-stacks}.}

With many processes, many threads and several CPUs, the relationships
among these structures must be recorded as well
(\cref{fig:l05-structures}):
\begin{itemize}
  \item the user-level library tracks all user-level threads of its process
    and the kernel-level threads on which it can run them;
  \item the PCB tracks the kernel-level threads that execute on behalf of
    the process;
  \item each kernel-level thread points to the PCB of its process, so that
    the kernel finds the address space when it runs the thread;
  \item a data structure for each CPU records, among other things, the
    thread that currently runs on it, and a running kernel-level thread
    points to its CPU.
\end{itemize}

\begin{figure}[htbp]
  \centering
  % Thread-related data structures at scale: user-level thread structures in
% the library, PCBs holding the hard process state, kernel-level threads
% holding the light process state, and CPU structures.
% Usage: % Thread-related data structures at scale: user-level thread structures in
% the library, PCBs holding the hard process state, kernel-level threads
% holding the light process state, and CPU structures.
% Usage: % Thread-related data structures at scale: user-level thread structures in
% the library, PCBs holding the hard process state, kernel-level threads
% holding the light process state, and CPU structures.
% Usage: \input{figures/l05-structures}   (width ~118 mm)
\begin{tikzpicture}[x=1mm, y=1mm, font=\scriptsize\sffamily,
  >={Stealth[length=1.5mm]},
  ult/.style={draw, fill=white, minimum width=10mm, minimum height=5mm,
              inner sep=0.5pt, align=center},
  klt/.style={draw, fill=white, minimum width=16mm, minimum height=11mm,
              inner sep=0.5pt, align=center},
  pcb/.style={draw, fill=ln-box, minimum height=14mm, inner sep=0.5pt,
              align=center},
  cpu/.style={draw, fill=ln-accent!22, minimum width=12mm, minimum height=5mm,
              inner sep=0.5pt},
  lab/.style={font=\scriptsize\sffamily, text=ln-muted},
  map/.style={densely dashed, ln-muted}]
  % user/kernel boundary
  \draw[dashed, ln-rule] (0,36) -- (118,36);
  \node[lab, anchor=south west] at (0,36.5) {\iflangja{ユーザ}{user}};
  \node[lab, anchor=north west] at (0,35.5) {\iflangja{カーネル}{kernel}};
  % ---- process P: library with user-level threads
  \draw[rounded corners=2pt, fill=ln-box] (22,39) rectangle (86,56);
  \node[lab, anchor=north west] at (22.5,55.8)
    {\iflangja{P のユーザレベルスレッドライブラリ}{user-level thread library of P}};
  \foreach \x/\n in {30/1, 44/2, 58/3, 72/4}
    \node[ult] (u\n) at (\x,45) {ULT\,\n};
  % ---- process P: PCB and kernel-level threads
  \node[pcb, minimum width=19mm] (pcbp) at (9.5,24)
    {\iflangja{P の PCB}{PCB of P}\\[1pt]
     \iflangja{ハード状態}{hard state}\\
     \iflangja{（アドレス}{(address}\\
     \iflangja{対応付け）}{mappings)}};
  \foreach \x/\n in {32/1, 54/2, 76/3}
    \node[klt] (k\n) at (\x,24)
      {\textbf{KLT\,\n}\\\iflangja{ライト状態}{light state}\\
       \iflangja{レジスタ等}{regs, stack}};
  \draw[->] (pcbp.east) -- (k1.west);
  \draw[->] (k1.east) -- (k2.west);
  \draw[->] (k2.east) -- (k3.west);
  % user-level threads currently mapped onto kernel-level threads
  \draw[map] (u1.south) -- (k1.north);
  \draw[map] (u3.south) -- (k2.north);
  \draw[map] (u4.south) -- (k3.north);
  % ---- process Q: single-threaded
  \draw[rounded corners=2pt, fill=ln-box] (93,39) rectangle (118,56);
  \node[align=center] at (105.5,47.5)
    {\iflangja{プロセス Q}{process Q}\\\iflangja{（シングル\\スレッド）}{(single-\\threaded)}};
  \node[pcb, minimum width=11mm] (pcbq) at (95.5,24)
    {\iflangja{Q の}{PCB}\\\iflangja{PCB}{of Q}};
  \node[klt, minimum width=14mm] (kq) at (111,24)
    {\textbf{KLT}\\\iflangja{ライト状態}{light state}\\
     \iflangja{レジスタ等}{regs, stack}};
  \draw[->] (pcbq.east) -- (kq.west);
  % ---- CPUs
  \node[cpu] (c0) at (32,5) {CPU\,0};
  \node[cpu] (c1) at (76,5) {CPU\,1};
  \node[cpu] (c2) at (111,5) {CPU\,2};
  \draw[<->] (k1.south) -- (c0.north);
  \draw[<->] (k3.south) -- (c1.north);
  \draw[<->] (kq.south) -- (c2.north);
  \node[lab, anchor=north] at (54,18) {\iflangja{実行中でない}{not running}};
\end{tikzpicture}
   (width ~118 mm)
\begin{tikzpicture}[x=1mm, y=1mm, font=\scriptsize\sffamily,
  >={Stealth[length=1.5mm]},
  ult/.style={draw, fill=white, minimum width=10mm, minimum height=5mm,
              inner sep=0.5pt, align=center},
  klt/.style={draw, fill=white, minimum width=16mm, minimum height=11mm,
              inner sep=0.5pt, align=center},
  pcb/.style={draw, fill=ln-box, minimum height=14mm, inner sep=0.5pt,
              align=center},
  cpu/.style={draw, fill=ln-accent!22, minimum width=12mm, minimum height=5mm,
              inner sep=0.5pt},
  lab/.style={font=\scriptsize\sffamily, text=ln-muted},
  map/.style={densely dashed, ln-muted}]
  % user/kernel boundary
  \draw[dashed, ln-rule] (0,36) -- (118,36);
  \node[lab, anchor=south west] at (0,36.5) {\iflangja{ユーザ}{user}};
  \node[lab, anchor=north west] at (0,35.5) {\iflangja{カーネル}{kernel}};
  % ---- process P: library with user-level threads
  \draw[rounded corners=2pt, fill=ln-box] (22,39) rectangle (86,56);
  \node[lab, anchor=north west] at (22.5,55.8)
    {\iflangja{P のユーザレベルスレッドライブラリ}{user-level thread library of P}};
  \foreach \x/\n in {30/1, 44/2, 58/3, 72/4}
    \node[ult] (u\n) at (\x,45) {ULT\,\n};
  % ---- process P: PCB and kernel-level threads
  \node[pcb, minimum width=19mm] (pcbp) at (9.5,24)
    {\iflangja{P の PCB}{PCB of P}\\[1pt]
     \iflangja{ハード状態}{hard state}\\
     \iflangja{（アドレス}{(address}\\
     \iflangja{対応付け）}{mappings)}};
  \foreach \x/\n in {32/1, 54/2, 76/3}
    \node[klt] (k\n) at (\x,24)
      {\textbf{KLT\,\n}\\\iflangja{ライト状態}{light state}\\
       \iflangja{レジスタ等}{regs, stack}};
  \draw[->] (pcbp.east) -- (k1.west);
  \draw[->] (k1.east) -- (k2.west);
  \draw[->] (k2.east) -- (k3.west);
  % user-level threads currently mapped onto kernel-level threads
  \draw[map] (u1.south) -- (k1.north);
  \draw[map] (u3.south) -- (k2.north);
  \draw[map] (u4.south) -- (k3.north);
  % ---- process Q: single-threaded
  \draw[rounded corners=2pt, fill=ln-box] (93,39) rectangle (118,56);
  \node[align=center] at (105.5,47.5)
    {\iflangja{プロセス Q}{process Q}\\\iflangja{（シングル\\スレッド）}{(single-\\threaded)}};
  \node[pcb, minimum width=11mm] (pcbq) at (95.5,24)
    {\iflangja{Q の}{PCB}\\\iflangja{PCB}{of Q}};
  \node[klt, minimum width=14mm] (kq) at (111,24)
    {\textbf{KLT}\\\iflangja{ライト状態}{light state}\\
     \iflangja{レジスタ等}{regs, stack}};
  \draw[->] (pcbq.east) -- (kq.west);
  % ---- CPUs
  \node[cpu] (c0) at (32,5) {CPU\,0};
  \node[cpu] (c1) at (76,5) {CPU\,1};
  \node[cpu] (c2) at (111,5) {CPU\,2};
  \draw[<->] (k1.south) -- (c0.north);
  \draw[<->] (k3.south) -- (c1.north);
  \draw[<->] (kq.south) -- (c2.north);
  \node[lab, anchor=north] at (54,18) {\iflangja{実行中でない}{not running}};
\end{tikzpicture}
   (width ~118 mm)
\begin{tikzpicture}[x=1mm, y=1mm, font=\scriptsize\sffamily,
  >={Stealth[length=1.5mm]},
  ult/.style={draw, fill=white, minimum width=10mm, minimum height=5mm,
              inner sep=0.5pt, align=center},
  klt/.style={draw, fill=white, minimum width=16mm, minimum height=11mm,
              inner sep=0.5pt, align=center},
  pcb/.style={draw, fill=ln-box, minimum height=14mm, inner sep=0.5pt,
              align=center},
  cpu/.style={draw, fill=ln-accent!22, minimum width=12mm, minimum height=5mm,
              inner sep=0.5pt},
  lab/.style={font=\scriptsize\sffamily, text=ln-muted},
  map/.style={densely dashed, ln-muted}]
  % user/kernel boundary
  \draw[dashed, ln-rule] (0,36) -- (118,36);
  \node[lab, anchor=south west] at (0,36.5) {\iflangja{ユーザ}{user}};
  \node[lab, anchor=north west] at (0,35.5) {\iflangja{カーネル}{kernel}};
  % ---- process P: library with user-level threads
  \draw[rounded corners=2pt, fill=ln-box] (22,39) rectangle (86,56);
  \node[lab, anchor=north west] at (22.5,55.8)
    {\iflangja{P のユーザレベルスレッドライブラリ}{user-level thread library of P}};
  \foreach \x/\n in {30/1, 44/2, 58/3, 72/4}
    \node[ult] (u\n) at (\x,45) {ULT\,\n};
  % ---- process P: PCB and kernel-level threads
  \node[pcb, minimum width=19mm] (pcbp) at (9.5,24)
    {\iflangja{P の PCB}{PCB of P}\\[1pt]
     \iflangja{ハード状態}{hard state}\\
     \iflangja{（アドレス}{(address}\\
     \iflangja{対応付け）}{mappings)}};
  \foreach \x/\n in {32/1, 54/2, 76/3}
    \node[klt] (k\n) at (\x,24)
      {\textbf{KLT\,\n}\\\iflangja{ライト状態}{light state}\\
       \iflangja{レジスタ等}{regs, stack}};
  \draw[->] (pcbp.east) -- (k1.west);
  \draw[->] (k1.east) -- (k2.west);
  \draw[->] (k2.east) -- (k3.west);
  % user-level threads currently mapped onto kernel-level threads
  \draw[map] (u1.south) -- (k1.north);
  \draw[map] (u3.south) -- (k2.north);
  \draw[map] (u4.south) -- (k3.north);
  % ---- process Q: single-threaded
  \draw[rounded corners=2pt, fill=ln-box] (93,39) rectangle (118,56);
  \node[align=center] at (105.5,47.5)
    {\iflangja{プロセス Q}{process Q}\\\iflangja{（シングル\\スレッド）}{(single-\\threaded)}};
  \node[pcb, minimum width=11mm] (pcbq) at (95.5,24)
    {\iflangja{Q の}{PCB}\\\iflangja{PCB}{of Q}};
  \node[klt, minimum width=14mm] (kq) at (111,24)
    {\textbf{KLT}\\\iflangja{ライト状態}{light state}\\
     \iflangja{レジスタ等}{regs, stack}};
  \draw[->] (pcbq.east) -- (kq.west);
  % ---- CPUs
  \node[cpu] (c0) at (32,5) {CPU\,0};
  \node[cpu] (c1) at (76,5) {CPU\,1};
  \node[cpu] (c2) at (111,5) {CPU\,2};
  \draw[<->] (k1.south) -- (c0.north);
  \draw[<->] (k3.south) -- (c1.north);
  \draw[<->] (kq.south) -- (c2.north);
  \node[lab, anchor=north] at (54,18) {\iflangja{実行中でない}{not running}};
\end{tikzpicture}

  \caption{Thread data structures at scale.  The PCB of a process keeps the
    state shared by all its threads and a list of its kernel-level threads
    (KLTs); each KLT keeps its own state and, while it runs, is linked to a
    CPU.  Dashed lines show which user-level threads (ULTs) the library has
    currently placed on which KLTs.  Back pointers from KLTs to PCBs are
    omitted.}
  \label{fig:l05-structures}
\end{figure}

\subsection{Hard and light process state}

When the kernel switches between two kernel-level threads of the same
process, the virtual address mappings remain valid, and only the state
that belongs to the individual thread has to be saved and restored.\tip{A
field is hard state if it is valid for every thread of the process, light
state if it is valid only for the thread currently executing.}  Some
of that state concerns the user-level threads: which signals are masked
and the arguments of a system call in progress belong to the user-level
thread that currently runs on a kernel-level thread, not to the whole
process.  The state of a process thus divides into two parts.

\needspace{9\baselineskip}
\begin{definition}[Hard and light process state]
  The \term{hard process state}[ハードプロセス状態] is the part of the
  state of a process that is relevant to all user-level threads executing
  within it, such as the virtual address mappings.  The
  \term{light process state}[ライトプロセス状態] is the part that is
  relevant only to the subset of user-level threads currently associated
  with a particular kernel-level thread, such as the signal mask and the
  system call arguments.\caution{The two terms come from the Solaris
    literature and from the lecture; textbooks usually speak instead of
    per-process and per-thread state.}
\end{definition}

\subsection{Why several data structures}

A single PCB and a set of smaller linked structures lead to different
costs (\cref{tab:l05-rationale}).  The single PCB is a large contiguous
structure, private to every execution context even where its contents are
shared; it is saved and restored in full on every context switch, and any
change of state updates it.  Smaller structures that refer to one another
through pointers keep shared state once, let a context switch save and
restore only what changes, and let the user-level library update the part
of the state it manages without involving the kernel.  The split improves
scalability, performance and flexibility, and reduces overhead.

\begin{table}[htbp]
  \centering
  \caption{A single PCB compared with several smaller data structures.}
  \label{tab:l05-rationale}
  \small
  \begin{tabular}{@{}>{\raggedright\arraybackslash}p{0.19\linewidth}>{\raggedright\arraybackslash}p{0.37\linewidth}>{\raggedright\arraybackslash}p{0.37\linewidth}@{}}
    \toprule
    & Single PCB & Several data structures \\
    \midrule
    Structure      & one large, contiguous structure
                   & small structures linked by pointers \\
    Sharing        & private to each execution context
                   & shared state kept once \\
    Context switch & saved and restored in full
                   & only the part that changes \\
    Updates        & every change modifies the PCB
                   & the library updates only its part \\
    \midrule
    Consequences   & poor scalability, high overhead, lower performance,
                     little flexibility
                   & better scalability, lower overhead, higher
                     performance, more flexibility \\
    \bottomrule
  \end{tabular}
\end{table}

\section{Threads in SunOS~5.0}
\label{sec:l05-solaris}

SunOS~5.0, the kernel of Solaris~2, was designed for shared-memory
multiprocessors, and the kernel itself is multithreaded.\source{P2L4,
SunOS~5.0 threading model, user-level thread data structures, kernel-level
data structures; Eykholt et al.\ \cite{eykholt1992}; Stein and Shah
\cite{stein1992}.}  Its threading model (\cref{fig:l05-sunos}) has the
following elements:
\begin{itemize}
  \item processes may be single-threaded or multithreaded; the user-level
    library supports both the many-to-many and the one-to-one model,
    including bound threads, and a process may run all its user-level
    threads on a single lightweight process;
  \item every kernel-level thread that executes user-level threads is
    associated with a lightweight process;
  \item other kernel-level threads have no lightweight process; they
    perform work of the kernel itself;
  \item the kernel-level scheduler schedules the kernel-level threads onto
    the CPUs.
\end{itemize}

\begin{definition}[Lightweight process]
  In SunOS~5.0 a \term{lightweight process}[軽量プロセス]%
  \index{LWP|see{lightweight process}} (LWP) is the kernel data structure
  associated with a kernel-level thread that executes user-level threads;
  it holds the light process state of that execution context.  To the
  user-level library an LWP represents a virtual CPU onto which it
  schedules user-level threads.
\end{definition}

\begin{figure}[htbp]
  \centering
  % The SunOS 5.0 threading model: user-level threads, lightweight processes,
% kernel-level threads, the kernel-level scheduler and the CPUs.
% Usage: % The SunOS 5.0 threading model: user-level threads, lightweight processes,
% kernel-level threads, the kernel-level scheduler and the CPUs.
% Usage: % The SunOS 5.0 threading model: user-level threads, lightweight processes,
% kernel-level threads, the kernel-level scheduler and the CPUs.
% Usage: \input{figures/l05-sunos}   (width ~118 mm)
\begin{tikzpicture}[x=1mm, y=1mm, font=\scriptsize\sffamily,
  ut/.style={draw, circle, fill=white, minimum size=3.6mm, inner sep=0pt},
  lwp/.style={draw, rounded corners=1pt, fill=ln-accent!15, minimum width=6.5mm,
              minimum height=3.6mm, inner sep=0pt, font=\tiny\sffamily},
  kt/.style={draw, rectangle, fill=ln-box, minimum size=3.6mm, inner sep=0pt},
  cpu/.style={draw, fill=ln-accent!22, minimum width=12mm, minimum height=4.5mm,
              inner sep=0pt},
  proc/.style={draw, dashed, rounded corners=2pt, ln-muted},
  lab/.style={font=\scriptsize\sffamily, text=ln-muted},
  ttl/.style={font=\scriptsize\sffamily\bfseries, anchor=base}]
  % user/kernel boundary runs through the LWPs
  \draw[dashed, ln-rule] (0,22) -- (118,22);
  \node[lab, anchor=south west] at (0,22.5) {\iflangja{ユーザ}{user}};
  \node[lab, anchor=north west] at (0,21.5) {\iflangja{カーネル}{kernel}};
  % ---- P1: many-to-many running on a single LWP
  \draw[proc] (12,17) rectangle (36,39);
  \node[ttl] at (24,40.5) {\iflangja{P1：LWP 1 個}{P1: one LWP}};
  \node[lwp] (l1) at (24,22) {LWP};
  \foreach \x in {17,24,31} { \node[ut] (u) at (\x,33) {}; \draw (u) -- (l1.north); }
  \node[kt] (k1) at (24,12) {};
  \draw (l1) -- (k1);
  % ---- P2: many-to-many with a bound thread
  \draw[proc] (40,17) rectangle (82,39);
  \node[ttl] at (61,40.5) {\iflangja{P2：多対多}{P2: many-to-many}};
  \node[lwp] (l2a) at (49,22) {LWP};
  \node[lwp] (l2b) at (60,22) {LWP};
  \node[ut] (u2a) at (45,33) {};
  \node[ut] (u2b) at (54.5,33) {};
  \node[ut] (u2c) at (64,33) {};
  \draw (u2a) -- (l2a.north);
  \draw (u2b) -- (l2a.north);
  \draw (u2b) -- (l2b.north);
  \draw (u2c) -- (l2b.north);
  \node[kt] (k2a) at (49,12) {};
  \node[kt] (k2b) at (60,12) {};
  \draw (l2a) -- (k2a);
  \draw (l2b) -- (k2b);
  \node[ut, thick, draw=ln-accent] (ub) at (75,33) {};
  \node[lwp, thick, draw=ln-accent] (lb) at (75,22) {LWP};
  \node[kt, thick, draw=ln-accent] (kb) at (75,12) {};
  \draw[very thick, ln-accent] (ub) -- (lb) -- (kb);
  \node[lab, text=ln-accent, anchor=south] at (75,35) {\iflangja{バウンド}{bound}};
  % ---- P3: one-to-one
  \draw[proc] (86,17) rectangle (106,39);
  \node[ttl] at (96,40.5) {\iflangja{P3：一対一}{P3: one-to-one}};
  \foreach \x/\n in {91/a, 101/b} {
    \node[ut] (u) at (\x,33) {};
    \node[lwp] (l) at (\x,22) {LWP};
    \node[kt] (k3\n) at (\x,12) {};
    \draw (u) -- (l) -- (k3\n);
  }
  % ---- kernel threads without LWP
  \node[kt] (kk1) at (111,12) {};
  \node[kt] (kk2) at (116,12) {};
  % ---- scheduler and CPUs
  \draw[rounded corners=1pt, fill=ln-box] (12,2) rectangle (118,7);
  \node at (65,4.5) {\iflangja{カーネルレベルスケジューラ}{kernel-level scheduler}};
  \foreach \k in {k1, k2a, k2b, kb, k3a, k3b, kk1, kk2}
    \draw (\k.south) -- (\k.south |- 0,7);
  \foreach \x/\n in {35/0, 65/1, 95/2} {
    \node[cpu] (c) at (\x,-6) {CPU\,\n};
    \draw (c.north) -- (\x,2);
  }
  % ---- legend
  \node[ut] at (14,-15) {};
  \node[lab, anchor=west] at (16.5,-15) {\iflangja{ユーザレベルスレッド}{user-level thread}};
  \node[lwp] at (48,-15) {LWP};
  \node[lab, anchor=west] at (52,-15) {\iflangja{軽量プロセス}{lightweight process}};
  \node[kt] at (84,-15) {};
  \node[lab, anchor=west] at (86.5,-15) {\iflangja{カーネルレベルスレッド}{kernel-level thread}};
\end{tikzpicture}
   (width ~118 mm)
\begin{tikzpicture}[x=1mm, y=1mm, font=\scriptsize\sffamily,
  ut/.style={draw, circle, fill=white, minimum size=3.6mm, inner sep=0pt},
  lwp/.style={draw, rounded corners=1pt, fill=ln-accent!15, minimum width=6.5mm,
              minimum height=3.6mm, inner sep=0pt, font=\tiny\sffamily},
  kt/.style={draw, rectangle, fill=ln-box, minimum size=3.6mm, inner sep=0pt},
  cpu/.style={draw, fill=ln-accent!22, minimum width=12mm, minimum height=4.5mm,
              inner sep=0pt},
  proc/.style={draw, dashed, rounded corners=2pt, ln-muted},
  lab/.style={font=\scriptsize\sffamily, text=ln-muted},
  ttl/.style={font=\scriptsize\sffamily\bfseries, anchor=base}]
  % user/kernel boundary runs through the LWPs
  \draw[dashed, ln-rule] (0,22) -- (118,22);
  \node[lab, anchor=south west] at (0,22.5) {\iflangja{ユーザ}{user}};
  \node[lab, anchor=north west] at (0,21.5) {\iflangja{カーネル}{kernel}};
  % ---- P1: many-to-many running on a single LWP
  \draw[proc] (12,17) rectangle (36,39);
  \node[ttl] at (24,40.5) {\iflangja{P1：LWP 1 個}{P1: one LWP}};
  \node[lwp] (l1) at (24,22) {LWP};
  \foreach \x in {17,24,31} { \node[ut] (u) at (\x,33) {}; \draw (u) -- (l1.north); }
  \node[kt] (k1) at (24,12) {};
  \draw (l1) -- (k1);
  % ---- P2: many-to-many with a bound thread
  \draw[proc] (40,17) rectangle (82,39);
  \node[ttl] at (61,40.5) {\iflangja{P2：多対多}{P2: many-to-many}};
  \node[lwp] (l2a) at (49,22) {LWP};
  \node[lwp] (l2b) at (60,22) {LWP};
  \node[ut] (u2a) at (45,33) {};
  \node[ut] (u2b) at (54.5,33) {};
  \node[ut] (u2c) at (64,33) {};
  \draw (u2a) -- (l2a.north);
  \draw (u2b) -- (l2a.north);
  \draw (u2b) -- (l2b.north);
  \draw (u2c) -- (l2b.north);
  \node[kt] (k2a) at (49,12) {};
  \node[kt] (k2b) at (60,12) {};
  \draw (l2a) -- (k2a);
  \draw (l2b) -- (k2b);
  \node[ut, thick, draw=ln-accent] (ub) at (75,33) {};
  \node[lwp, thick, draw=ln-accent] (lb) at (75,22) {LWP};
  \node[kt, thick, draw=ln-accent] (kb) at (75,12) {};
  \draw[very thick, ln-accent] (ub) -- (lb) -- (kb);
  \node[lab, text=ln-accent, anchor=south] at (75,35) {\iflangja{バウンド}{bound}};
  % ---- P3: one-to-one
  \draw[proc] (86,17) rectangle (106,39);
  \node[ttl] at (96,40.5) {\iflangja{P3：一対一}{P3: one-to-one}};
  \foreach \x/\n in {91/a, 101/b} {
    \node[ut] (u) at (\x,33) {};
    \node[lwp] (l) at (\x,22) {LWP};
    \node[kt] (k3\n) at (\x,12) {};
    \draw (u) -- (l) -- (k3\n);
  }
  % ---- kernel threads without LWP
  \node[kt] (kk1) at (111,12) {};
  \node[kt] (kk2) at (116,12) {};
  % ---- scheduler and CPUs
  \draw[rounded corners=1pt, fill=ln-box] (12,2) rectangle (118,7);
  \node at (65,4.5) {\iflangja{カーネルレベルスケジューラ}{kernel-level scheduler}};
  \foreach \k in {k1, k2a, k2b, kb, k3a, k3b, kk1, kk2}
    \draw (\k.south) -- (\k.south |- 0,7);
  \foreach \x/\n in {35/0, 65/1, 95/2} {
    \node[cpu] (c) at (\x,-6) {CPU\,\n};
    \draw (c.north) -- (\x,2);
  }
  % ---- legend
  \node[ut] at (14,-15) {};
  \node[lab, anchor=west] at (16.5,-15) {\iflangja{ユーザレベルスレッド}{user-level thread}};
  \node[lwp] at (48,-15) {LWP};
  \node[lab, anchor=west] at (52,-15) {\iflangja{軽量プロセス}{lightweight process}};
  \node[kt] at (84,-15) {};
  \node[lab, anchor=west] at (86.5,-15) {\iflangja{カーネルレベルスレッド}{kernel-level thread}};
\end{tikzpicture}
   (width ~118 mm)
\begin{tikzpicture}[x=1mm, y=1mm, font=\scriptsize\sffamily,
  ut/.style={draw, circle, fill=white, minimum size=3.6mm, inner sep=0pt},
  lwp/.style={draw, rounded corners=1pt, fill=ln-accent!15, minimum width=6.5mm,
              minimum height=3.6mm, inner sep=0pt, font=\tiny\sffamily},
  kt/.style={draw, rectangle, fill=ln-box, minimum size=3.6mm, inner sep=0pt},
  cpu/.style={draw, fill=ln-accent!22, minimum width=12mm, minimum height=4.5mm,
              inner sep=0pt},
  proc/.style={draw, dashed, rounded corners=2pt, ln-muted},
  lab/.style={font=\scriptsize\sffamily, text=ln-muted},
  ttl/.style={font=\scriptsize\sffamily\bfseries, anchor=base}]
  % user/kernel boundary runs through the LWPs
  \draw[dashed, ln-rule] (0,22) -- (118,22);
  \node[lab, anchor=south west] at (0,22.5) {\iflangja{ユーザ}{user}};
  \node[lab, anchor=north west] at (0,21.5) {\iflangja{カーネル}{kernel}};
  % ---- P1: many-to-many running on a single LWP
  \draw[proc] (12,17) rectangle (36,39);
  \node[ttl] at (24,40.5) {\iflangja{P1：LWP 1 個}{P1: one LWP}};
  \node[lwp] (l1) at (24,22) {LWP};
  \foreach \x in {17,24,31} { \node[ut] (u) at (\x,33) {}; \draw (u) -- (l1.north); }
  \node[kt] (k1) at (24,12) {};
  \draw (l1) -- (k1);
  % ---- P2: many-to-many with a bound thread
  \draw[proc] (40,17) rectangle (82,39);
  \node[ttl] at (61,40.5) {\iflangja{P2：多対多}{P2: many-to-many}};
  \node[lwp] (l2a) at (49,22) {LWP};
  \node[lwp] (l2b) at (60,22) {LWP};
  \node[ut] (u2a) at (45,33) {};
  \node[ut] (u2b) at (54.5,33) {};
  \node[ut] (u2c) at (64,33) {};
  \draw (u2a) -- (l2a.north);
  \draw (u2b) -- (l2a.north);
  \draw (u2b) -- (l2b.north);
  \draw (u2c) -- (l2b.north);
  \node[kt] (k2a) at (49,12) {};
  \node[kt] (k2b) at (60,12) {};
  \draw (l2a) -- (k2a);
  \draw (l2b) -- (k2b);
  \node[ut, thick, draw=ln-accent] (ub) at (75,33) {};
  \node[lwp, thick, draw=ln-accent] (lb) at (75,22) {LWP};
  \node[kt, thick, draw=ln-accent] (kb) at (75,12) {};
  \draw[very thick, ln-accent] (ub) -- (lb) -- (kb);
  \node[lab, text=ln-accent, anchor=south] at (75,35) {\iflangja{バウンド}{bound}};
  % ---- P3: one-to-one
  \draw[proc] (86,17) rectangle (106,39);
  \node[ttl] at (96,40.5) {\iflangja{P3：一対一}{P3: one-to-one}};
  \foreach \x/\n in {91/a, 101/b} {
    \node[ut] (u) at (\x,33) {};
    \node[lwp] (l) at (\x,22) {LWP};
    \node[kt] (k3\n) at (\x,12) {};
    \draw (u) -- (l) -- (k3\n);
  }
  % ---- kernel threads without LWP
  \node[kt] (kk1) at (111,12) {};
  \node[kt] (kk2) at (116,12) {};
  % ---- scheduler and CPUs
  \draw[rounded corners=1pt, fill=ln-box] (12,2) rectangle (118,7);
  \node at (65,4.5) {\iflangja{カーネルレベルスケジューラ}{kernel-level scheduler}};
  \foreach \k in {k1, k2a, k2b, kb, k3a, k3b, kk1, kk2}
    \draw (\k.south) -- (\k.south |- 0,7);
  \foreach \x/\n in {35/0, 65/1, 95/2} {
    \node[cpu] (c) at (\x,-6) {CPU\,\n};
    \draw (c.north) -- (\x,2);
  }
  % ---- legend
  \node[ut] at (14,-15) {};
  \node[lab, anchor=west] at (16.5,-15) {\iflangja{ユーザレベルスレッド}{user-level thread}};
  \node[lwp] at (48,-15) {LWP};
  \node[lab, anchor=west] at (52,-15) {\iflangja{軽量プロセス}{lightweight process}};
  \node[kt] at (84,-15) {};
  \node[lab, anchor=west] at (86.5,-15) {\iflangja{カーネルレベルスレッド}{kernel-level thread}};
\end{tikzpicture}

  \caption{The SunOS~5.0 threading model.  Each kernel-level thread that
    runs user-level threads has an LWP; the two kernel-level threads on the
    right perform kernel work and have none.  P1 runs all its user-level
    threads on a single LWP.  The kernel-level scheduler places
    kernel-level threads on CPUs.}
  \label{fig:l05-sunos}
\end{figure}

\subsection{User-level thread data structures}

When the library of Stein and Shah creates a thread, it returns a thread
identifier.  The identifier is not a pointer to the thread data structure
but an index into a table of pointers to the structures.  The indirection
protects against stale identifiers: if a thread has exited and the memory
of its structure has been reused, a direct pointer would silently refer to
unrelated data, whereas the table entry can record the state of the thread
so that the library returns a meaningful error.

The thread data structure contains the execution context (the registers,
including the stack pointer), the signal mask, the priority, the
\emph{thread-local storage}, and the stack.  Thread-local storage holds the
variables that the thread functions declare as private to each thread;
they are known at compile time, so the compiler can reserve space for them
in every thread.  The stack size is set by a library default or by the
application.  Since the size of all parts is known when the thread is
created, the library can allocate the structures of successive threads
next to one another in contiguous memory, which gives good locality.

Contiguous allocation has a danger.  The library does not control how far
a stack grows, and the operating system does not know about user-level
threads at all.  A thread whose stack overflows therefore writes into the
data structure of the neighbouring thread, and the damage shows only when
that other thread runs and fails---far from the cause.

\begin{definition}[Red zone]
  A \term{red zone}[レッドゾーン] is a region of virtual address space
  between a thread's stack and the neighbouring thread's data structure
  that is not mapped to memory, so that any access to it causes a
  fault.\caution{An unrelated meaning: in the x86-64 ABI the \emph{red zone}
    is the \SI{128}{\byte} below the stack pointer that a leaf function may
    use without adjusting it.}
\end{definition}

A stack that grows into its red zone causes a fault at once, in the thread
whose stack overflowed, which makes the error easy to diagnose
(\cref{fig:l05-ult-layout}).

\begin{figure}[htbp]
  \centering
  % User-level thread data structures allocated contiguously, with a red zone
% below each stack, and the table of pointers indexed by thread identifier.
% Usage: % User-level thread data structures allocated contiguously, with a red zone
% below each stack, and the table of pointers indexed by thread identifier.
% Usage: % User-level thread data structures allocated contiguously, with a red zone
% below each stack, and the table of pointers indexed by thread identifier.
% Usage: \input{figures/l05-ult-layout}   (width ~116 mm)
\begin{tikzpicture}[x=1mm, y=1mm, font=\scriptsize\sffamily,
  >={Stealth[length=1.5mm]},
  lab/.style={font=\scriptsize\sffamily, text=ln-muted, align=center},
  cell/.style={draw, fill=white, minimum width=8mm, minimum height=5mm,
               inner sep=0pt}]
  \node[lab] at (3,4.5) {\ldots};
  \foreach \o/\t in {8/1, 60/2} {
    \draw[fill=ln-caution!25] (\o,0) rectangle ++(5,9);
    \draw[fill=white] (\o+5,0) rectangle ++(30,9);
    \node at (\o+20,6.2) {\iflangja{T\t のスタック}{stack of T\t}};
    \draw[->, ln-accent, thick] (\o+29,2.5) -- (\o+11,2.5);
    \draw[fill=ln-box] (\o+35,0) rectangle ++(6,9);
    \node at (\o+38,4.5) {TLS};
    \draw[fill=ln-accent!22] (\o+41,0) rectangle ++(11,9);
    \node[align=center] at (\o+46.5,4.5) {T\t\\\iflangja{構造体}{struct}};
    \node[lab, text=ln-caution, anchor=north] at (\o+2.5,-0.8)
      {\iflangja{レッドゾーン}{red zone}};
  }
  \node[lab] at (115,4.5) {\ldots};
  % address axis
  \draw[->, ln-muted] (8,-7.5) -- (112,-7.5);
  \node[lab, anchor=north west] at (8,-8) {\iflangja{下位アドレス}{lower addresses}};
  \node[lab, anchor=north east] at (112,-8) {\iflangja{上位アドレス}{higher addresses}};
  \node[lab, anchor=north] at (60,-8)
    {\iflangja{スタックは下位アドレスへ伸びる}{stacks grow toward lower addresses}};
  % table of pointers indexed by thread id
  \node[lab, anchor=east] at (33,24) {\iflangja{ポインタの表}{table of pointers}};
  \foreach \i/\x in {1/38, 2/46, 3/54} {
    \node[cell] (t\i) at (\x,24) {};
    \node[lab, anchor=south] at (\x,26.5) {\i};
  }
  \node[cell] at (62,24) {\ldots};
  \node[lab, anchor=west] at (67,24) {\iflangja{（添字 = スレッド ID）}{(index = thread ID)}};
  \fill (t1) circle (0.6);
  \fill (t2) circle (0.6);
  \draw[->] (t1.center) -- (54.5,9);
  \draw[->] (t2.center) .. controls (46,15) and (100,16) .. (106.5,9);
\end{tikzpicture}
   (width ~116 mm)
\begin{tikzpicture}[x=1mm, y=1mm, font=\scriptsize\sffamily,
  >={Stealth[length=1.5mm]},
  lab/.style={font=\scriptsize\sffamily, text=ln-muted, align=center},
  cell/.style={draw, fill=white, minimum width=8mm, minimum height=5mm,
               inner sep=0pt}]
  \node[lab] at (3,4.5) {\ldots};
  \foreach \o/\t in {8/1, 60/2} {
    \draw[fill=ln-caution!25] (\o,0) rectangle ++(5,9);
    \draw[fill=white] (\o+5,0) rectangle ++(30,9);
    \node at (\o+20,6.2) {\iflangja{T\t のスタック}{stack of T\t}};
    \draw[->, ln-accent, thick] (\o+29,2.5) -- (\o+11,2.5);
    \draw[fill=ln-box] (\o+35,0) rectangle ++(6,9);
    \node at (\o+38,4.5) {TLS};
    \draw[fill=ln-accent!22] (\o+41,0) rectangle ++(11,9);
    \node[align=center] at (\o+46.5,4.5) {T\t\\\iflangja{構造体}{struct}};
    \node[lab, text=ln-caution, anchor=north] at (\o+2.5,-0.8)
      {\iflangja{レッドゾーン}{red zone}};
  }
  \node[lab] at (115,4.5) {\ldots};
  % address axis
  \draw[->, ln-muted] (8,-7.5) -- (112,-7.5);
  \node[lab, anchor=north west] at (8,-8) {\iflangja{下位アドレス}{lower addresses}};
  \node[lab, anchor=north east] at (112,-8) {\iflangja{上位アドレス}{higher addresses}};
  \node[lab, anchor=north] at (60,-8)
    {\iflangja{スタックは下位アドレスへ伸びる}{stacks grow toward lower addresses}};
  % table of pointers indexed by thread id
  \node[lab, anchor=east] at (33,24) {\iflangja{ポインタの表}{table of pointers}};
  \foreach \i/\x in {1/38, 2/46, 3/54} {
    \node[cell] (t\i) at (\x,24) {};
    \node[lab, anchor=south] at (\x,26.5) {\i};
  }
  \node[cell] at (62,24) {\ldots};
  \node[lab, anchor=west] at (67,24) {\iflangja{（添字 = スレッド ID）}{(index = thread ID)}};
  \fill (t1) circle (0.6);
  \fill (t2) circle (0.6);
  \draw[->] (t1.center) -- (54.5,9);
  \draw[->] (t2.center) .. controls (46,15) and (100,16) .. (106.5,9);
\end{tikzpicture}
   (width ~116 mm)
\begin{tikzpicture}[x=1mm, y=1mm, font=\scriptsize\sffamily,
  >={Stealth[length=1.5mm]},
  lab/.style={font=\scriptsize\sffamily, text=ln-muted, align=center},
  cell/.style={draw, fill=white, minimum width=8mm, minimum height=5mm,
               inner sep=0pt}]
  \node[lab] at (3,4.5) {\ldots};
  \foreach \o/\t in {8/1, 60/2} {
    \draw[fill=ln-caution!25] (\o,0) rectangle ++(5,9);
    \draw[fill=white] (\o+5,0) rectangle ++(30,9);
    \node at (\o+20,6.2) {\iflangja{T\t のスタック}{stack of T\t}};
    \draw[->, ln-accent, thick] (\o+29,2.5) -- (\o+11,2.5);
    \draw[fill=ln-box] (\o+35,0) rectangle ++(6,9);
    \node at (\o+38,4.5) {TLS};
    \draw[fill=ln-accent!22] (\o+41,0) rectangle ++(11,9);
    \node[align=center] at (\o+46.5,4.5) {T\t\\\iflangja{構造体}{struct}};
    \node[lab, text=ln-caution, anchor=north] at (\o+2.5,-0.8)
      {\iflangja{レッドゾーン}{red zone}};
  }
  \node[lab] at (115,4.5) {\ldots};
  % address axis
  \draw[->, ln-muted] (8,-7.5) -- (112,-7.5);
  \node[lab, anchor=north west] at (8,-8) {\iflangja{下位アドレス}{lower addresses}};
  \node[lab, anchor=north east] at (112,-8) {\iflangja{上位アドレス}{higher addresses}};
  \node[lab, anchor=north] at (60,-8)
    {\iflangja{スタックは下位アドレスへ伸びる}{stacks grow toward lower addresses}};
  % table of pointers indexed by thread id
  \node[lab, anchor=east] at (33,24) {\iflangja{ポインタの表}{table of pointers}};
  \foreach \i/\x in {1/38, 2/46, 3/54} {
    \node[cell] (t\i) at (\x,24) {};
    \node[lab, anchor=south] at (\x,26.5) {\i};
  }
  \node[cell] at (62,24) {\ldots};
  \node[lab, anchor=west] at (67,24) {\iflangja{（添字 = スレッド ID）}{(index = thread ID)}};
  \fill (t1) circle (0.6);
  \fill (t2) circle (0.6);
  \draw[->] (t1.center) -- (54.5,9);
  \draw[->] (t2.center) .. controls (46,15) and (100,16) .. (106.5,9);
\end{tikzpicture}

  \caption{User-level thread data structures allocated contiguously.  The
    thread identifier indexes a table of pointers to the structures.  An
    overflow of the stack of T2 runs into the red zone before it can reach
    the structure of T1.}
  \label{fig:l05-ult-layout}
\end{figure}

\needspace{13\baselineskip}
\begin{example}
  A library gives each thread a block of \SI{64}{\kibi\byte}: a
  \SI{4}{\kibi\byte} red zone at the bottom, a \SI{56}{\kibi\byte} stack,
  and \SI{4}{\kibi\byte} at the top for the thread structure and its
  thread-local storage.  Measured in \si{\kibi\byte} from the start of the
  region, block $k$ occupies $[64k, 64k+64)$.  Thread~3 thus has its red
  zone at $[192, 196)$, its stack at $[196, 252)$ and its structure at
  $[252, 256)$, while the structure of thread~2 lies at $[188, 192)$.  If
  thread~3 needs more than \SI{56}{\kibi\byte} of stack, its next push
  lands in $[192, 196)$ and faults; without the red zone it would overwrite
  the structure of thread~2.  For 1024 threads the region spans
  \SI{64}{\mebi\byte} of virtual address space, of which the
  \SI{4}{\mebi\byte} of red zones use no physical memory.\margin{On x86,
    NPTL keeps the same layout: the thread descriptor sits at the top of the
    stack block, and glibc puts one guard page below the stack---%
    \SI{4}{\kibi\byte} by default, \texttt{pthread\_attr\_setguardsize}.}
\end{example}

\subsection{Kernel-level data structures}

The kernel of SunOS~5.0 keeps four kinds of structures for processes and
threads (\cref{tab:l05-kernel-structs}).  The LWP resembles the data
structure of a user-level thread, but it is visible to the kernel.  Its
contents are needed only while the process runs, so the LWP can be swapped
out together with the process.  The kernel-level thread structure holds
information that the kernel needs even when the process is not running,
such as scheduling information, and it is never swapped out.  The CPU
structure identifies the thread currently running on each
CPU.\margin{On SPARC processors, which have spare registers, SunOS keeps
the pointer to the current thread in a dedicated register.}

\begin{table}[htbp]
  \centering
  \caption{Kernel data structures for processes and threads in SunOS~5.0.}
  \label{tab:l05-kernel-structs}
  \small
  \begin{tabular}{@{}>{\raggedright\arraybackslash}p{0.22\linewidth}>{\raggedright\arraybackslash}p{0.71\linewidth}@{}}
    \toprule
    Structure & Contents \\
    \midrule
    Process & list of kernel-level threads, virtual address space, user
      credentials, signal handlers \\
    LWP & user-level registers, system call arguments, resource usage
      information, signal mask; swappable \\
    Kernel-level thread & kernel-level registers, stack pointer,
      scheduling information, pointers to the associated LWP, process and
      CPU; not swappable \\
    CPU & current thread, list of kernel-level threads, dispatching and
      interrupt-handling information \\
    \bottomrule
  \end{tabular}
\end{table}

\section{Interaction between the library and the kernel}
\label{sec:l05-interaction}

The user-level library does not know what happens in the kernel, and the
kernel does not know what happens in the library.\source{P2L4, basic
thread management interactions, thread management visibility and design;
Stein and Shah \cite{stein1992}; Silberschatz et al.\ ch.~4.}  System
calls and special signals let the two levels coordinate.  The application
tells the library how many of its threads should run at the same
time---its \emph{concurrency level}, set with
\texttt{pthread\_setconcurrency}---and the library asks the kernel for
kernel-level threads through system calls, requesting more when it needs
them; the kernel can notify the library of events that affect its threads,
such as a kernel-level thread blocking.

\needspace{6\baselineskip}
\begin{example}
  A process has six user-level threads, of which at most three need to run
  at the same time.  \Cref{tab:l05-interaction} traces how the library and
  the kernel adjust the number of kernel-level threads.\margin{On Linux
    \texttt{pthread\_setconcurrency} has no effect: with a one-to-one
    library there is nothing to choose, and POSIX lets an implementation
    ignore the hint.}
\end{example}

\begin{table}[htbp]
  \centering
  \caption{Coordination between a user-level library and the kernel for a
    process with six user-level threads (U1--U6).}
  \label{tab:l05-interaction}
  \small
  \begin{tabular}{@{}c>{\raggedright\arraybackslash}p{0.66\linewidth}>{\raggedright\arraybackslash}p{0.18\linewidth}@{}}
    \toprule
    Step & Event & KLTs \\
    \midrule
    1 & The library asks the kernel for three kernel-level threads; the
        kernel creates K1--K3. & K1--K3 \\
    2 & U1--U3 run on K1--K3; U4--U6 wait on a user-level condition. &
        3 running \\
    3 & U1 and U2 issue blocking I/O requests, and K1 and K2 block in
        the kernel.  Meanwhile U4 becomes ready. & 2 blocked, 1 running \\
    4 & U3 also issues a blocking request, so no kernel-level thread is
        left to run U4 and the whole process would stall.  When K3 blocks,
        the kernel therefore notifies the library---in SunOS~5.0 with the
        signal \texttt{SIGWAITING}; the library requests another
        kernel-level thread, the kernel creates K4, and U4 runs on it. &
        3 blocked, 1 running \\
    5 & The I/O completes and U1--U3 continue.  K4 then stays idle for a
        long time, and the library destroys it. & K1--K3 \\
    \bottomrule
  \end{tabular}
\end{table}

\subsection{Visibility of thread management}

Each level sees only part of the picture.
\begin{description}
  \item[The kernel] sees the kernel-level threads, the CPUs and the
    decisions of the kernel-level scheduler.
  \item[The user-level library] sees the user-level threads and the
    kernel-level threads assigned to its process.
\end{description}
The library may bind a user-level thread to a kernel-level thread, as for
the bound threads of Lesson~3, and the kernel may pin a kernel-level
thread to a CPU; neither decision is visible at the other
level.\margin{In POSIX a bound thread is \texttt{PTHREAD\_SCOPE\_SYSTEM},
the only scope NPTL supports; pinning a kernel-level thread to CPUs is the
cache affinity of Lesson~7.}

The lack of visibility causes problems.  Suppose a user-level thread
acquires a mutex implemented by the library, and the kernel then preempts
the kernel-level thread on which it runs, for instance because its
timeslice has expired.  The other user-level threads that need the mutex
cannot proceed, even when the kernel runs the kernel-level threads on which
they are scheduled: nothing can change until the kernel runs the preempted
kernel-level thread again, and the kernel, which does not see the mutex,
has no reason to do so soon.  In the many-to-many model, user-level
scheduling decisions and changes to the mapping of user-level threads onto
kernel-level threads are visible only to the library.  The one-to-one model
avoids many of these problems, because the kernel sees every
thread.\sidenote{The general remedy for the many-to-many model is to make
the kernel report every event that changes the number of running threads.
Anderson et al.\ \cite{anderson1992} called such an upcall a
\emph{scheduler activation}: the kernel hands the library a fresh activation
whenever a thread blocks, unblocks or is preempted, so that the library can
reschedule at once; \texttt{SIGWAITING} is a narrow form of the same idea.
NetBSD and FreeBSD implemented activations and later removed them---like the
many-to-many model itself, they cost more than they returned.}

Being ordinary code in the address space of the process, the library's
scheduler runs only when execution enters the library:
\begin{itemize}
  \item a user-level thread explicitly yields;
  \item a timer set by the library expires;
  \item a user-level thread calls a library function, such as locking or
    unlocking a mutex, that can change which threads are runnable;
  \item a blocked user-level thread becomes runnable;
  \item a signal arrives from the kernel, for example on behalf of a timer.
\end{itemize}

\section{Implementation issues}
\label{sec:l05-issues}

Multiprocessors raise further problems for the library, and the cost of
creating threads suggests reusing them.\source{P2L4, issues on multiple
CPUs, synchronization-related issues, destroying threads; Eykholt et al.\
\cite{eykholt1992}; OSTEP ch.~28.}

\subsection{User-level threads on multiple CPUs}

On a multiprocessor the library runs on several kernel-level threads at
the same time, and a decision taken on one CPU can require an action on
another.  Consider user-level threads T1, T2 and T3 with priorities
$\text{T3} > \text{T2} > \text{T1}$, which the library runs on two
kernel-level threads, K1 and K2, on two CPUs (\cref{fig:l05-multi-cpu}).
T2 holds mutex \texttt{m} on CPU~1.  T3 runs on CPU~2 and tries to lock
\texttt{m}; it blocks, and the library schedules T1 on K2.  When T2
unlocks \texttt{m}, the library, executing on CPU~1, finds that T3 is
runnable again and has a higher priority than T1, which is running.  T1
must be preempted, but the library code on CPU~1 cannot modify the
registers of CPU~2.  It sends a signal to K2; the signal makes K2 execute
the library on CPU~2, which preempts T1 and schedules T3.

\begin{figure}[htbp]
  \centering
  % Priority scheduling of user-level threads on two CPUs: when T2 unlocks the
% mutex, the library must signal the other kernel-level thread to preempt T1.
% Usage: % Priority scheduling of user-level threads on two CPUs: when T2 unlocks the
% mutex, the library must signal the other kernel-level thread to preempt T1.
% Usage: % Priority scheduling of user-level threads on two CPUs: when T2 unlocks the
% mutex, the library must signal the other kernel-level thread to preempt T1.
% Usage: \input{figures/l05-multi-cpu}   (width ~116 mm)
\begin{tikzpicture}[x=1mm, y=1mm, font=\footnotesize\sffamily,
  lab/.style={font=\scriptsize\sffamily, align=center},
  row/.style={font=\scriptsize\sffamily, anchor=east, align=right},
  >={Stealth[length=1.6mm]}]
  % \lfseg{x1}{x2}{y}{fill}{text}
  \def\lfseg#1#2#3#4#5{%
    \draw[fill=#4] (#1,#3-2.5) rectangle (#2,#3+2.5);
    \node[lab] at ({(#1+#2)/2},#3) {#5};}
  % event lines
  \draw[densely dashed, ln-muted] (44,24) -- (44,-2);
  \draw[densely dashed, ln-muted] (74,24) -- (74,-2);
  \node[lab, anchor=south, fill=white, inner sep=1pt] at (44,24)
    {\iflangja{T3 が m でブロック}{T3 blocks on m}};
  \node[lab, anchor=south, fill=white, inner sep=1pt] at (74,24)
    {\iflangja{T2 が m を解放}{T2 unlocks m}};
  % CPU 1
  \node[row] at (24,15) {CPU 1\\(KLT K1)};
  \lfseg{26}{74}{15}{ln-box}{\iflangja{T2（m を保持）}{T2 (holds m)}}
  \lfseg{74}{116}{15}{white}{T2}
  % CPU 2
  \node[row] at (24,4) {CPU 2\\(KLT K2)};
  \lfseg{26}{44}{4}{white}{T3}
  \lfseg{44}{48}{4}{black!35}{}
  \lfseg{48}{79}{4}{white}{T1}
  \lfseg{79}{83}{4}{black!35}{}
  \lfseg{83}{116}{4}{ln-box}{\iflangja{T3（m を保持）}{T3 (holds m)}}
  % signal from K1 to K2
  \draw[->, thick, ln-caution] (74,12.5) -- node[lab, right, pos=0.45, text=ln-caution]
    {\iflangja{シグナル}{signal}} (79,6.5);
  % legend
  \draw[fill=black!35] (26,-9) rectangle (30,-6);
  \node[lab, anchor=west] at (31,-7.5) {\iflangja{ライブラリのスケジューラ}{library scheduler}};
  \draw[fill=ln-box] (64,-9) rectangle (68,-6);
  \node[lab, anchor=west] at (69,-7.5) {\iflangja{m を保持}{m held}};
  \draw[->] (92,-7.5) -- node[lab, above] {\iflangja{時間}{time}} (116,-7.5);
\end{tikzpicture}
   (width ~116 mm)
\begin{tikzpicture}[x=1mm, y=1mm, font=\footnotesize\sffamily,
  lab/.style={font=\scriptsize\sffamily, align=center},
  row/.style={font=\scriptsize\sffamily, anchor=east, align=right},
  >={Stealth[length=1.6mm]}]
  % \lfseg{x1}{x2}{y}{fill}{text}
  \def\lfseg#1#2#3#4#5{%
    \draw[fill=#4] (#1,#3-2.5) rectangle (#2,#3+2.5);
    \node[lab] at ({(#1+#2)/2},#3) {#5};}
  % event lines
  \draw[densely dashed, ln-muted] (44,24) -- (44,-2);
  \draw[densely dashed, ln-muted] (74,24) -- (74,-2);
  \node[lab, anchor=south, fill=white, inner sep=1pt] at (44,24)
    {\iflangja{T3 が m でブロック}{T3 blocks on m}};
  \node[lab, anchor=south, fill=white, inner sep=1pt] at (74,24)
    {\iflangja{T2 が m を解放}{T2 unlocks m}};
  % CPU 1
  \node[row] at (24,15) {CPU 1\\(KLT K1)};
  \lfseg{26}{74}{15}{ln-box}{\iflangja{T2（m を保持）}{T2 (holds m)}}
  \lfseg{74}{116}{15}{white}{T2}
  % CPU 2
  \node[row] at (24,4) {CPU 2\\(KLT K2)};
  \lfseg{26}{44}{4}{white}{T3}
  \lfseg{44}{48}{4}{black!35}{}
  \lfseg{48}{79}{4}{white}{T1}
  \lfseg{79}{83}{4}{black!35}{}
  \lfseg{83}{116}{4}{ln-box}{\iflangja{T3（m を保持）}{T3 (holds m)}}
  % signal from K1 to K2
  \draw[->, thick, ln-caution] (74,12.5) -- node[lab, right, pos=0.45, text=ln-caution]
    {\iflangja{シグナル}{signal}} (79,6.5);
  % legend
  \draw[fill=black!35] (26,-9) rectangle (30,-6);
  \node[lab, anchor=west] at (31,-7.5) {\iflangja{ライブラリのスケジューラ}{library scheduler}};
  \draw[fill=ln-box] (64,-9) rectangle (68,-6);
  \node[lab, anchor=west] at (69,-7.5) {\iflangja{m を保持}{m held}};
  \draw[->] (92,-7.5) -- node[lab, above] {\iflangja{時間}{time}} (116,-7.5);
\end{tikzpicture}
   (width ~116 mm)
\begin{tikzpicture}[x=1mm, y=1mm, font=\footnotesize\sffamily,
  lab/.style={font=\scriptsize\sffamily, align=center},
  row/.style={font=\scriptsize\sffamily, anchor=east, align=right},
  >={Stealth[length=1.6mm]}]
  % \lfseg{x1}{x2}{y}{fill}{text}
  \def\lfseg#1#2#3#4#5{%
    \draw[fill=#4] (#1,#3-2.5) rectangle (#2,#3+2.5);
    \node[lab] at ({(#1+#2)/2},#3) {#5};}
  % event lines
  \draw[densely dashed, ln-muted] (44,24) -- (44,-2);
  \draw[densely dashed, ln-muted] (74,24) -- (74,-2);
  \node[lab, anchor=south, fill=white, inner sep=1pt] at (44,24)
    {\iflangja{T3 が m でブロック}{T3 blocks on m}};
  \node[lab, anchor=south, fill=white, inner sep=1pt] at (74,24)
    {\iflangja{T2 が m を解放}{T2 unlocks m}};
  % CPU 1
  \node[row] at (24,15) {CPU 1\\(KLT K1)};
  \lfseg{26}{74}{15}{ln-box}{\iflangja{T2（m を保持）}{T2 (holds m)}}
  \lfseg{74}{116}{15}{white}{T2}
  % CPU 2
  \node[row] at (24,4) {CPU 2\\(KLT K2)};
  \lfseg{26}{44}{4}{white}{T3}
  \lfseg{44}{48}{4}{black!35}{}
  \lfseg{48}{79}{4}{white}{T1}
  \lfseg{79}{83}{4}{black!35}{}
  \lfseg{83}{116}{4}{ln-box}{\iflangja{T3（m を保持）}{T3 (holds m)}}
  % signal from K1 to K2
  \draw[->, thick, ln-caution] (74,12.5) -- node[lab, right, pos=0.45, text=ln-caution]
    {\iflangja{シグナル}{signal}} (79,6.5);
  % legend
  \draw[fill=black!35] (26,-9) rectangle (30,-6);
  \node[lab, anchor=west] at (31,-7.5) {\iflangja{ライブラリのスケジューラ}{library scheduler}};
  \draw[fill=ln-box] (64,-9) rectangle (68,-6);
  \node[lab, anchor=west] at (69,-7.5) {\iflangja{m を保持}{m held}};
  \draw[->] (92,-7.5) -- node[lab, above] {\iflangja{時間}{time}} (116,-7.5);
\end{tikzpicture}

  \caption{User-level threads with priorities
    $\text{T3} > \text{T2} > \text{T1}$ on two CPUs.  When T2 unlocks
    \texttt{m}, the library on CPU~1 signals K2 so that the library on
    CPU~2 preempts T1 in favour of T3.}
  \label{fig:l05-multi-cpu}
\end{figure}

\subsection{Adaptive mutexes}

A multiprocessor also changes the best way to wait for a mutex.  Suppose
one thread holds a mutex on CPU~1 when another thread, running on CPU~2,
tries to lock it.  Normally the second thread is placed on the wait queue
of the mutex and blocks, which costs a context switch now and another one
when it is woken up.  If the
critical section is short and its owner is running on another CPU, the
owner is likely to release the mutex before these context switches would
complete, and it is cheaper for the second thread to \emph{spin}: to stay
on its CPU and repeatedly check whether the mutex is free.  If the owner is not running,
spinning is useless, since the mutex cannot be released before the owner
is scheduled again.

\begin{definition}[Adaptive mutex]
  An \term{adaptive mutex}[適応型ミューテックス] is a mutex whose lock
  operation spins while the owner of the mutex is running on another CPU,
  and blocks the calling thread otherwise.
\end{definition}

Adaptive mutexes make sense only on multiprocessors---on a single CPU the
owner cannot be running while another thread spins---and only for short
critical sections; long critical sections are better served by blocking.
To decide between spinning and blocking, the lock operation must find out
whether the owner is running, so the mutex data structure records its
owner.  Lesson~10 returns to spinning as the basis of the
spinlock.\margin{OSTEP calls a lock that first spins for a while and then
sleeps a \emph{two-phase lock}; the mutex of the Linux kernel spins only
while the owner is running, exactly as here.}

\subsection{Destroying threads}

Creating a thread allocates and initializes its data structures and its
stack, and destroying it frees them.  Because applications create and
destroy threads frequently, reusing these structures pays off.  When a
thread exits, its structures are not freed immediately; the thread is
placed on a list called the \emph{death row}.  A \term{reaper thread}%
[リーパスレッド] runs periodically, destroys the threads on the death row
and frees their data structures.  If a request to create a thread arrives
before the reaper has run, the structure and stack of a thread on the
death row are reused, which saves the cost of allocating and initializing
them.\margin{glibc does the same for NPTL threads: a freed stack stays in a
cache of up to \SI{40}{\mebi\byte} (tunable
\texttt{glibc.pthread.stack\_cache\_size}) and is handed to the next
\texttt{pthread\_create}.}

\needspace{11\baselineskip}
\begin{keypoint}
  The kernel and the user-level library each see only their own threads.
  System calls and signals let them coordinate, but decisions such as
  preempting a lock owner or preempting a thread on another CPU need
  explicit notification across the levels.  Several CPUs also change how a
  thread waits---an adaptive mutex spins only while the owner runs---and
  threads that exit are kept for reuse instead of being freed at once.
\end{keypoint}

\section{Interrupts and signals}
\label{sec:l05-int-sig}

Two kinds of events can divert a thread from its normal execution:
interrupts and signals.\source{P2L4, interrupts and signals intro; OSTEP
ch.~6; Kerrisk \cite{kerrisk2010} ch.~20.}  Signals were introduced in
Lesson~1 as notifications that the operating system passes to an
application.  Interrupts are the corresponding notifications that the
hardware delivers to the CPU.

\begin{definition}[Interrupt]
  An \term{interrupt}[割り込み] is an event generated externally to the
  CPU, by a component such as an I/O device, a timer or another CPU, that
  notifies the CPU that some condition requires attention.
\end{definition}

Which interrupts can occur is determined by the physical platform, and
interrupts appear asynchronously, independently of what the CPU is
executing.  Signals, in contrast, are events triggered by the
CPU and the software running on it; which signals exist is determined by
the operating system, and they can appear synchronously or asynchronously.
The two mechanisms nevertheless have much in common
(\cref{tab:l05-int-vs-sig}): both have a unique identifier, defined by the
hardware or by the operating system; both can be masked, that is, disabled
or suppressed; and when enabled, both trigger a handler.  Interrupt
handlers are set by the operating system for the entire system, signal
handlers by each process for itself.

\begin{table}[htbp]
  \centering
  \caption{Interrupts and signals.}
  \label{tab:l05-int-vs-sig}
  \small
  \begin{tabular}{@{}>{\raggedright\arraybackslash}p{0.22\linewidth}>{\raggedright\arraybackslash}p{0.345\linewidth}>{\raggedright\arraybackslash}p{0.345\linewidth}@{}}
    \toprule
    & Interrupt & Signal \\
    \midrule
    Generated by & devices, timers, other CPUs & the CPU and the software
      running on it \\
    Defined by   & the hardware platform & the operating system \\
    Occurrence   & asynchronous & synchronous or asynchronous \\
    Masked       & per CPU & per execution context \\
    Handler set  & by the operating system, system-wide & by each process
      \\
    \bottomrule
  \end{tabular}
\end{table}

\section{Handling interrupts and signals}
\label{sec:l05-handling}

When a device needs attention, it sends an interrupt to the CPU over the
interconnect that links them, such as PCI Express.\source{P2L4, interrupt
handling, signal handling; OSTEP ch.~6, 36; Kerrisk ch.~20--21.}  Early
systems used a dedicated physical line for every interrupt; modern devices
send the interrupt as a message over the interconnect itself, and the
content of the message identifies the interrupt and thereby the device.
An interrupt raised by sending such a message instead of asserting a
dedicated line is a
\term{message signaled interrupt}[メッセージシグナル割り込み]%
\index{MSI|see{message signaled interrupt}} (MSI).

\begin{definition}[Interrupt handler]
  An \term{interrupt handler}[割り込みハンドラ] is the routine that the
  operating system executes in response to a particular interrupt.  The
  \term{interrupt handler table}[割り込みハンドラテーブル] maps each
  interrupt number to the start address of its handler.
\end{definition}

The CPU uses the interrupt number to look up the interrupt handler table,
sets its program counter to the start address found there, and executes
the handler (\cref{fig:l05-int-sig}a).  Which interrupts occur is specific
to the hardware; how they are handled is specific to the operating
system.\margin{The table is also called the interrupt vector table; on x86
it is the interrupt descriptor table.}

\begin{figure}[htbp]
  \centering
  % Handling an interrupt (system-wide handler table) and a signal (per-process
% handler table).
% Usage: % Handling an interrupt (system-wide handler table) and a signal (per-process
% handler table).
% Usage: % Handling an interrupt (system-wide handler table) and a signal (per-process
% handler table).
% Usage: \input{figures/l05-int-sig}   (width ~118 mm)
\begin{tikzpicture}[x=1mm, y=1mm, font=\scriptsize\sffamily,
  >={Stealth[length=1.5mm]},
  box/.style={draw, fill=white, minimum height=9mm, inner sep=1pt, align=center},
  act/.style={draw, fill=ln-accent!22, minimum height=9mm, inner sep=1pt,
              align=center},
  lab/.style={font=\scriptsize\sffamily, align=center},
  ttl/.style={font=\footnotesize\sffamily\bfseries, anchor=west},
  tcap/.style={font=\scriptsize\sffamily, text=ln-muted, anchor=south}]
  % \lstable{y-top}{row1a}{row1b}{row3a}{row3b}: three-row handler table at x=52..88
  \def\lstable#1#2#3#4#5{%
    \draw[fill=ln-box] (52,#1) rectangle (88,#1-15);
    \draw (52,#1-5) -- (88,#1-5);
    \draw (52,#1-10) -- (88,#1-10);
    \draw (65,#1) -- (65,#1-15);
    \node at (58.5,#1-2.5) {#2};  \node at (76.5,#1-2.5) {#3};
    \node at (58.5,#1-7.5) {\ldots}; \node at (76.5,#1-7.5) {\ldots};
    \node at (58.5,#1-12.5) {#4}; \node at (76.5,#1-12.5) {#5};}
  % ---- (a) interrupt
  \node[ttl] at (0,33) {\iflangja{(a) 割り込み}{(a) interrupt}};
  \node[box, minimum width=16mm] (dev) at (8,20) {\iflangja{デバイス}{device}};
  \node[act, minimum width=12mm] (cpu) at (38,20) {CPU};
  \draw[->, thick] (dev.east) -- node[lab, above] {MSI} node[lab, below]
    {\iflangja{（PCIe）}{(PCIe)}} (cpu.west);
  \lstable{27.5}{INT-1}{\iflangja{ハンドラ 1 の番地}{addr.\ handler 1}}%
    {INT-N}{\iflangja{ハンドラ N の番地}{addr.\ handler N}}
  \node[tcap] at (70,27.8) {\iflangja{割り込みハンドラテーブル（システム全体）}{interrupt handler table (system-wide)}};
  \draw[->] (cpu.east) -- node[lab, above] {N} (47,20) |- (52,15);
  \node[act, minimum width=20mm] (h1) at (107,15) {\iflangja{ハンドラ N}{handler N}};
  \draw[->, thick] (88,15) -- node[lab, above] {\iflangja{PC を}{set}}
    node[lab, below] {\iflangja{設定}{PC}} (h1.west);
  % ---- (b) signal
  \node[ttl] at (0,1) {\iflangja{(b) シグナル}{(b) signal}};
  \node[box, minimum width=16mm] (thr) at (8,-12) {\iflangja{P の\\スレッド}{thread\\of P}};
  \node[act, minimum width=12mm] (ker) at (38,-12) {\iflangja{カーネル}{kernel}};
  \draw[->, thick] (thr.east) -- node[lab, above] {\iflangja{不正な}{illegal}}
    node[lab, below] {\iflangja{アクセス}{access}} (ker.west);
  \lstable{-4.5}{SIG-1}{\iflangja{ハンドラ / 既定}{handler / default}}%
    {SIGSEGV}{\iflangja{ハンドラの番地}{addr.\ handler}}
  \node[tcap] at (70,-4.2) {\iflangja{P のシグナルハンドラテーブル}{signal handler table of P}};
  \draw[->] (ker.east) -- node[lab, above] {11} (47,-12) |- (52,-17);
  \node[act, minimum width=20mm] (h2) at (107,-17) {\iflangja{ハンドラ\\（スレッドの\\文脈で実行）}{handler\\(in thread's\\context)}};
  \draw[->, thick] (88,-17) -- node[lab, above] {\iflangja{PC を}{set}}
    node[lab, below] {\iflangja{設定}{PC}} (h2.west);
\end{tikzpicture}
   (width ~118 mm)
\begin{tikzpicture}[x=1mm, y=1mm, font=\scriptsize\sffamily,
  >={Stealth[length=1.5mm]},
  box/.style={draw, fill=white, minimum height=9mm, inner sep=1pt, align=center},
  act/.style={draw, fill=ln-accent!22, minimum height=9mm, inner sep=1pt,
              align=center},
  lab/.style={font=\scriptsize\sffamily, align=center},
  ttl/.style={font=\footnotesize\sffamily\bfseries, anchor=west},
  tcap/.style={font=\scriptsize\sffamily, text=ln-muted, anchor=south}]
  % \lstable{y-top}{row1a}{row1b}{row3a}{row3b}: three-row handler table at x=52..88
  \def\lstable#1#2#3#4#5{%
    \draw[fill=ln-box] (52,#1) rectangle (88,#1-15);
    \draw (52,#1-5) -- (88,#1-5);
    \draw (52,#1-10) -- (88,#1-10);
    \draw (65,#1) -- (65,#1-15);
    \node at (58.5,#1-2.5) {#2};  \node at (76.5,#1-2.5) {#3};
    \node at (58.5,#1-7.5) {\ldots}; \node at (76.5,#1-7.5) {\ldots};
    \node at (58.5,#1-12.5) {#4}; \node at (76.5,#1-12.5) {#5};}
  % ---- (a) interrupt
  \node[ttl] at (0,33) {\iflangja{(a) 割り込み}{(a) interrupt}};
  \node[box, minimum width=16mm] (dev) at (8,20) {\iflangja{デバイス}{device}};
  \node[act, minimum width=12mm] (cpu) at (38,20) {CPU};
  \draw[->, thick] (dev.east) -- node[lab, above] {MSI} node[lab, below]
    {\iflangja{（PCIe）}{(PCIe)}} (cpu.west);
  \lstable{27.5}{INT-1}{\iflangja{ハンドラ 1 の番地}{addr.\ handler 1}}%
    {INT-N}{\iflangja{ハンドラ N の番地}{addr.\ handler N}}
  \node[tcap] at (70,27.8) {\iflangja{割り込みハンドラテーブル（システム全体）}{interrupt handler table (system-wide)}};
  \draw[->] (cpu.east) -- node[lab, above] {N} (47,20) |- (52,15);
  \node[act, minimum width=20mm] (h1) at (107,15) {\iflangja{ハンドラ N}{handler N}};
  \draw[->, thick] (88,15) -- node[lab, above] {\iflangja{PC を}{set}}
    node[lab, below] {\iflangja{設定}{PC}} (h1.west);
  % ---- (b) signal
  \node[ttl] at (0,1) {\iflangja{(b) シグナル}{(b) signal}};
  \node[box, minimum width=16mm] (thr) at (8,-12) {\iflangja{P の\\スレッド}{thread\\of P}};
  \node[act, minimum width=12mm] (ker) at (38,-12) {\iflangja{カーネル}{kernel}};
  \draw[->, thick] (thr.east) -- node[lab, above] {\iflangja{不正な}{illegal}}
    node[lab, below] {\iflangja{アクセス}{access}} (ker.west);
  \lstable{-4.5}{SIG-1}{\iflangja{ハンドラ / 既定}{handler / default}}%
    {SIGSEGV}{\iflangja{ハンドラの番地}{addr.\ handler}}
  \node[tcap] at (70,-4.2) {\iflangja{P のシグナルハンドラテーブル}{signal handler table of P}};
  \draw[->] (ker.east) -- node[lab, above] {11} (47,-12) |- (52,-17);
  \node[act, minimum width=20mm] (h2) at (107,-17) {\iflangja{ハンドラ\\（スレッドの\\文脈で実行）}{handler\\(in thread's\\context)}};
  \draw[->, thick] (88,-17) -- node[lab, above] {\iflangja{PC を}{set}}
    node[lab, below] {\iflangja{設定}{PC}} (h2.west);
\end{tikzpicture}
   (width ~118 mm)
\begin{tikzpicture}[x=1mm, y=1mm, font=\scriptsize\sffamily,
  >={Stealth[length=1.5mm]},
  box/.style={draw, fill=white, minimum height=9mm, inner sep=1pt, align=center},
  act/.style={draw, fill=ln-accent!22, minimum height=9mm, inner sep=1pt,
              align=center},
  lab/.style={font=\scriptsize\sffamily, align=center},
  ttl/.style={font=\footnotesize\sffamily\bfseries, anchor=west},
  tcap/.style={font=\scriptsize\sffamily, text=ln-muted, anchor=south}]
  % \lstable{y-top}{row1a}{row1b}{row3a}{row3b}: three-row handler table at x=52..88
  \def\lstable#1#2#3#4#5{%
    \draw[fill=ln-box] (52,#1) rectangle (88,#1-15);
    \draw (52,#1-5) -- (88,#1-5);
    \draw (52,#1-10) -- (88,#1-10);
    \draw (65,#1) -- (65,#1-15);
    \node at (58.5,#1-2.5) {#2};  \node at (76.5,#1-2.5) {#3};
    \node at (58.5,#1-7.5) {\ldots}; \node at (76.5,#1-7.5) {\ldots};
    \node at (58.5,#1-12.5) {#4}; \node at (76.5,#1-12.5) {#5};}
  % ---- (a) interrupt
  \node[ttl] at (0,33) {\iflangja{(a) 割り込み}{(a) interrupt}};
  \node[box, minimum width=16mm] (dev) at (8,20) {\iflangja{デバイス}{device}};
  \node[act, minimum width=12mm] (cpu) at (38,20) {CPU};
  \draw[->, thick] (dev.east) -- node[lab, above] {MSI} node[lab, below]
    {\iflangja{（PCIe）}{(PCIe)}} (cpu.west);
  \lstable{27.5}{INT-1}{\iflangja{ハンドラ 1 の番地}{addr.\ handler 1}}%
    {INT-N}{\iflangja{ハンドラ N の番地}{addr.\ handler N}}
  \node[tcap] at (70,27.8) {\iflangja{割り込みハンドラテーブル（システム全体）}{interrupt handler table (system-wide)}};
  \draw[->] (cpu.east) -- node[lab, above] {N} (47,20) |- (52,15);
  \node[act, minimum width=20mm] (h1) at (107,15) {\iflangja{ハンドラ N}{handler N}};
  \draw[->, thick] (88,15) -- node[lab, above] {\iflangja{PC を}{set}}
    node[lab, below] {\iflangja{設定}{PC}} (h1.west);
  % ---- (b) signal
  \node[ttl] at (0,1) {\iflangja{(b) シグナル}{(b) signal}};
  \node[box, minimum width=16mm] (thr) at (8,-12) {\iflangja{P の\\スレッド}{thread\\of P}};
  \node[act, minimum width=12mm] (ker) at (38,-12) {\iflangja{カーネル}{kernel}};
  \draw[->, thick] (thr.east) -- node[lab, above] {\iflangja{不正な}{illegal}}
    node[lab, below] {\iflangja{アクセス}{access}} (ker.west);
  \lstable{-4.5}{SIG-1}{\iflangja{ハンドラ / 既定}{handler / default}}%
    {SIGSEGV}{\iflangja{ハンドラの番地}{addr.\ handler}}
  \node[tcap] at (70,-4.2) {\iflangja{P のシグナルハンドラテーブル}{signal handler table of P}};
  \draw[->] (ker.east) -- node[lab, above] {11} (47,-12) |- (52,-17);
  \node[act, minimum width=20mm] (h2) at (107,-17) {\iflangja{ハンドラ\\（スレッドの\\文脈で実行）}{handler\\(in thread's\\context)}};
  \draw[->, thick] (88,-17) -- node[lab, above] {\iflangja{PC を}{set}}
    node[lab, below] {\iflangja{設定}{PC}} (h2.west);
\end{tikzpicture}

  \caption{(a) An interrupt raised by a device selects a handler from the
    system-wide interrupt handler table.  (b) An illegal memory access
    makes the kernel raise \texttt{SIGSEGV}, which selects a handler from
    the signal handler table of the process.}
  \label{fig:l05-int-sig}
\end{figure}

Signals are generated by the operating system.  When a thread accesses
memory that is not allocated to its process, for example, the operating
system raises \texttt{SIGSEGV} (signal~11 on Linux).\margin{Linux defines
31 standard signals and a block of real-time ones; glibc reserves the first
two of those for NPTL, so \texttt{SIGRTMIN} is 34 and \texttt{SIGRTMAX} 64
(\texttt{signal(7)}).}  For each process it
keeps a \emph{signal handler table} that maps every signal number to the
action to be taken (\cref{fig:l05-int-sig}b).  The set of signals is
operating-system specific; POSIX defines a standard set.  Every signal has
a default action, which is one of the following: terminate the process,
ignore the signal, terminate the process and dump core, stop the process,
or continue a stopped process.

\begin{definition}[Signal handler]
  A \term{signal handler}[シグナルハンドラ] is a function that a process
  installs to be executed, in the context of the receiving thread, when a
  particular signal is delivered.
\end{definition}

A process installs handlers with \texttt{signal} or, more portably, with
\texttt{sigaction}.  Most signals can be caught in this way; a few cannot,
notably \texttt{SIGKILL} and \texttt{SIGSTOP}, whose default actions always
apply.

\begin{example}
  A program, using the headers \texttt{signal.h} and \texttt{unistd.h},
  that stops its work loop when a timer expires:\caution{A handler may
    assign only to \texttt{volatile sig\_atomic\_t} objects and call the
    functions listed in \texttt{signal-safety(7)}; \texttt{printf} and
    \texttt{malloc} are not among them.}
  \begin{lstlisting}[language=C, numbers=none]
static volatile sig_atomic_t expired = 0;

static void on_alarm(int signo) {
    (void)signo;
    expired = 1;            /* keep handlers minimal */
}

int main(void) {
    struct sigaction sa;
    sa.sa_handler = on_alarm;
    sigemptyset(&sa.sa_mask);
    sa.sa_flags = 0;
    sigaction(SIGALRM, &sa, NULL);
    alarm(2);               /* SIGALRM after 2 s */
    while (!expired)
        do_work();
    return 0;
}
\end{lstlisting}
\end{example}

Signals are classified by how they relate to the thread that receives
them.

\begin{definition}[Synchronous and asynchronous signals]
  A \term{synchronous signal}[同期シグナル] occurs in response to an
  action of a specific thread and is delivered to that thread, for example
  \texttt{SIGSEGV} on an access to protected memory or \texttt{SIGFPE} on a
  division by zero.  An
  \term{asynchronous signal}[非同期シグナル] occurs independently of what
  the receiving thread executes, for example \texttt{SIGKILL} sent by
  another process with \texttt{kill}, or \texttt{SIGALRM} on the
  expiration of a timer.
\end{definition}

A signal can also be sent explicitly to one thread, for example with
\texttt{pthread\_kill}; like a synchronous signal, it is then delivered to
that thread only.

\section{Disabling interrupts and signals}
\label{sec:l05-masks}

Interrupt and signal handlers execute in the context of the thread that
was interrupted: the program counter of that thread is set to the handler,
and the handler runs on its behalf, without the scheduler being
involved---a signal handler on the thread's own stack, an interrupt
handler on a kernel stack.\source{P2L4, why disable
interrupts or signals, more on signal masks, interrupts on multicore
systems, types of signals; Kerrisk ch.~20--22.}  This creates a problem
when the handler accesses shared state.  If the interrupted thread holds a
mutex \texttt{m} and the handler tries to lock \texttt{m}, the handler
blocks.  The owner of \texttt{m} is the interrupted thread itself, which
cannot continue until the handler returns: the thread deadlocks with its
own handler.

There are two ways to avoid this.  Handler code can be kept simple, never
locking a mutex; this is often too restrictive.  Or the code can control
when handlers may run, by disabling particular interrupts or signals while
it holds mutexes that their handlers need.

\begin{definition}[Interrupt mask and signal mask]
  An \term{interrupt mask}[割り込みマスク] or a
  \term{signal mask}[シグナルマスク] is a sequence of bits with one bit per
  interrupt or signal; the value of the bit indicates whether the interrupt
  or signal is enabled (1) or disabled (0).
\end{definition}

\needspace{7\baselineskip}
A thread that disables a signal around a critical section follows this
pattern:
\begin{lstlisting}[language=C, numbers=none]
clear_field_in_mask(mask);  // disable
lock(mutex) {
  ...                       // handler cannot run here
} // unlock
reset_field_in_mask(mask);  // enable
\end{lstlisting}
An interrupt or signal that occurs while it is disabled is not lost: it
remains \emph{pending}, and its handler runs as soon as the mask enables it
again.  If the same signal occurs several times while it is disabled, the
handler normally runs only once; generating a signal several times does
not guarantee that the handler runs several times (\cref{sec:l05-types}).
\caution{POSIX signal sets list the \emph{blocked} signals: a signal that is
a member of the set is disabled, the opposite of the bit convention
above.}

\subsection{Setting the masks}

Interrupt masks are kept per CPU.  If the mask of a CPU disables an
interrupt, the hardware interrupt-routing mechanism does not deliver that
interrupt to the CPU.  The kernel sets the mask by writing to the
hardware.  Signal masks are kept per execution context: the kernel does
not interrupt a kernel-level thread whose mask disables the signal, and
where user-level threads run on top of kernel-level threads the kernel
sees only the kernel-level masks, which \cref{sec:l05-thread-signals}
returns to.  A single-threaded process changes its signal mask with
\texttt{sigprocmask}, a thread of a multithreaded process with
\texttt{pthread\_sigmask}; both take one of the operations
\texttt{SIG\_BLOCK}, \texttt{SIG\_UNBLOCK} or \texttt{SIG\_SETMASK}.

\needspace{11\baselineskip}
\begin{example}
  The pattern above written with POSIX calls, for a handler of
  \texttt{SIGALRM} that locks the same mutex:
  \begin{lstlisting}[language=C, numbers=none]
sigset_t set, old;
sigemptyset(&set);
sigaddset(&set, SIGALRM);
pthread_sigmask(SIG_BLOCK, &set, &old);    // disable
pthread_mutex_lock(&m);
/* ... critical section ... */
pthread_mutex_unlock(&m);
pthread_sigmask(SIG_SETMASK, &old, NULL);  // restore
\end{lstlisting}
  \caution{\texttt{pthread\_mutex\_lock} is not async-signal-safe; a
    handler should rather set a flag, and every thread that may run it
    must block the signal.}
\end{example}

\subsection{Interrupts on multicore systems}

On a multicore system an interrupt can be directed to any CPU that has it
enabled.  Enabling a given interrupt on a single core only, so that only
that core handles it, avoids the overhead and the perturbation that
handling it would cause on all the other cores.

\needspace{5\baselineskip}
\begin{keypoint}
  Handlers run in the context of the interrupted thread.  Code that holds
  a mutex needed by a handler must disable the corresponding interrupt or
  signal while it holds the mutex; an event that occurs meanwhile stays
  pending until it is enabled again.
\end{keypoint}

\subsection{Types of signals}
\label{sec:l05-types}

Signals differ in what happens when several instances arrive before the
handler has run.
\begin{description}
  \item[One-shot signals] With a \term{one-shot signal}[ワンショットシグナル],
    $n$ pending instances are equivalent to one: the handler is guaranteed
    to run at least once, but not $n$ times.  On some systems the handler
    is also reset to the default action after it has run, so the process
    must re-install it explicitly to catch the signal again.
  \item[Real-time signals] With a
    \term{real-time signal}[リアルタイムシグナル], instances are queued:
    if the signal is raised $n$ times, the handler is called $n$ times.
    POSIX reserves the range \texttt{SIGRTMIN} to \texttt{SIGRTMAX} for
    them.
\end{description}

\begin{example}
  A thread blocks both \texttt{SIGUSR1}, a standard signal, and the
  real-time signal \texttt{SIGRTMIN}.  While they are blocked, each signal
  is raised four times; then the thread unblocks them.  The handler of
  \texttt{SIGUSR1} runs once, because the four instances collapse into one
  pending signal; the handler of \texttt{SIGRTMIN} runs four times.
\end{example}

\section{Interrupts as threads}
\label{sec:l05-int-threads}

SunOS~5.0 avoids the deadlock between a thread and its handler in a
different way: it handles interrupts as threads.\source{P2L4, interrupts
as threads, performance of threads as interrupts; Eykholt et al.\
\cite{eykholt1992}.}  If an interrupt handler runs as a thread of its own,
it can block on a mutex like any other thread; it waits until the owner,
possibly the interrupted thread, releases the mutex, and the scheduler
runs it afterwards.  Creating a new thread for every interrupt would be
too expensive, so interrupt threads are created and initialized in advance,
and whether one of them becomes a full thread is decided at run
time.  An interrupt switches onto the
stack of such a preallocated thread, while the interrupted thread stays
pinned beneath it and cannot run; a handler that does not block finishes
there, and the pinned thread resumes.  Only if the handler blocks on a
mutex does the interrupt thread become a complete thread, which the
scheduler treats like any other, and the pinned thread continues at once.

The same idea underlies a common way of dividing the work of interrupt
handling (\cref{fig:l05-top-bottom}).\margin{Of the Linux deferral
mechanisms, only work queues and threaded handlers may block.}

\begin{definition}[Top half and bottom half]
  The \term{top half}[トップハーフ] of interrupt handling runs immediately
  when the interrupt occurs; it does a minimal amount of processing and
  must be fast and non-blocking.  The \term{bottom half}[ボトムハーフ] is
  the deferred part; it can be of arbitrary complexity and, when a thread
  executes it, it is scheduled like any other thread and may block.
\end{definition}

\begin{figure}[htbp]
  \centering
  % Interrupt handling split into a top half, executed immediately, and a
% bottom half, executed later like a thread.
% Usage: % Interrupt handling split into a top half, executed immediately, and a
% bottom half, executed later like a thread.
% Usage: % Interrupt handling split into a top half, executed immediately, and a
% bottom half, executed later like a thread.
% Usage: \input{figures/l05-top-bottom}   (width ~116 mm)
\begin{tikzpicture}[x=1mm, y=1mm, font=\footnotesize\sffamily,
  lab/.style={font=\scriptsize\sffamily, align=center},
  row/.style={font=\scriptsize\sffamily, anchor=east},
  >={Stealth[length=1.6mm]}]
  \def\ltbseg#1#2#3#4#5{%
    \draw[fill=#4] (#1,#3-2.5) rectangle (#2,#3+2.5);
    \node[lab] at ({(#1+#2)/2},#3) {#5};}
  \node[row] at (14,10) {CPU};
  \ltbseg{16}{44}{10}{white}{\iflangja{スレッド T}{thread T}}
  \ltbseg{44}{61}{10}{ln-accent!22}{\iflangja{トップハーフ}{top half}}
  \ltbseg{61}{80}{10}{white}{\iflangja{スレッド T}{thread T}}
  \ltbseg{80}{104}{10}{ln-box}{\iflangja{ボトムハーフ}{bottom half}}
  \ltbseg{104}{116}{10}{white}{T}
  % interrupt arrival
  \draw[->, thick, ln-caution] (44,20) -- (44,12.5);
  \node[lab, anchor=south, text=ln-caution] at (44,20) {\iflangja{割り込み}{interrupt}};
  % top half defers the bottom half
  \draw[->, ln-accent] (52.5,7.5) .. controls (56,1) and (86,1) .. (92,7.5);
  \node[lab, text=ln-accent] at (72,4.8) {\iflangja{後で実行するよう登録}{deferred}};
  % annotations
  \node[lab, anchor=north] at (52.5,-1)
    {\iflangja{直ちに実行\\高速・ブロックしない}{runs immediately;\\fast, non-blocking}};
  \node[lab, anchor=north] at (92,-1)
    {\iflangja{スレッドとしてスケジュール\\ブロック可能}{scheduled like a thread;\\may block}};
  \draw[->] (92,20) -- node[lab, above] {\iflangja{時間}{time}} (116,20);
\end{tikzpicture}
   (width ~116 mm)
\begin{tikzpicture}[x=1mm, y=1mm, font=\footnotesize\sffamily,
  lab/.style={font=\scriptsize\sffamily, align=center},
  row/.style={font=\scriptsize\sffamily, anchor=east},
  >={Stealth[length=1.6mm]}]
  \def\ltbseg#1#2#3#4#5{%
    \draw[fill=#4] (#1,#3-2.5) rectangle (#2,#3+2.5);
    \node[lab] at ({(#1+#2)/2},#3) {#5};}
  \node[row] at (14,10) {CPU};
  \ltbseg{16}{44}{10}{white}{\iflangja{スレッド T}{thread T}}
  \ltbseg{44}{61}{10}{ln-accent!22}{\iflangja{トップハーフ}{top half}}
  \ltbseg{61}{80}{10}{white}{\iflangja{スレッド T}{thread T}}
  \ltbseg{80}{104}{10}{ln-box}{\iflangja{ボトムハーフ}{bottom half}}
  \ltbseg{104}{116}{10}{white}{T}
  % interrupt arrival
  \draw[->, thick, ln-caution] (44,20) -- (44,12.5);
  \node[lab, anchor=south, text=ln-caution] at (44,20) {\iflangja{割り込み}{interrupt}};
  % top half defers the bottom half
  \draw[->, ln-accent] (52.5,7.5) .. controls (56,1) and (86,1) .. (92,7.5);
  \node[lab, text=ln-accent] at (72,4.8) {\iflangja{後で実行するよう登録}{deferred}};
  % annotations
  \node[lab, anchor=north] at (52.5,-1)
    {\iflangja{直ちに実行\\高速・ブロックしない}{runs immediately;\\fast, non-blocking}};
  \node[lab, anchor=north] at (92,-1)
    {\iflangja{スレッドとしてスケジュール\\ブロック可能}{scheduled like a thread;\\may block}};
  \draw[->] (92,20) -- node[lab, above] {\iflangja{時間}{time}} (116,20);
\end{tikzpicture}
   (width ~116 mm)
\begin{tikzpicture}[x=1mm, y=1mm, font=\footnotesize\sffamily,
  lab/.style={font=\scriptsize\sffamily, align=center},
  row/.style={font=\scriptsize\sffamily, anchor=east},
  >={Stealth[length=1.6mm]}]
  \def\ltbseg#1#2#3#4#5{%
    \draw[fill=#4] (#1,#3-2.5) rectangle (#2,#3+2.5);
    \node[lab] at ({(#1+#2)/2},#3) {#5};}
  \node[row] at (14,10) {CPU};
  \ltbseg{16}{44}{10}{white}{\iflangja{スレッド T}{thread T}}
  \ltbseg{44}{61}{10}{ln-accent!22}{\iflangja{トップハーフ}{top half}}
  \ltbseg{61}{80}{10}{white}{\iflangja{スレッド T}{thread T}}
  \ltbseg{80}{104}{10}{ln-box}{\iflangja{ボトムハーフ}{bottom half}}
  \ltbseg{104}{116}{10}{white}{T}
  % interrupt arrival
  \draw[->, thick, ln-caution] (44,20) -- (44,12.5);
  \node[lab, anchor=south, text=ln-caution] at (44,20) {\iflangja{割り込み}{interrupt}};
  % top half defers the bottom half
  \draw[->, ln-accent] (52.5,7.5) .. controls (56,1) and (86,1) .. (92,7.5);
  \node[lab, text=ln-accent] at (72,4.8) {\iflangja{後で実行するよう登録}{deferred}};
  % annotations
  \node[lab, anchor=north] at (52.5,-1)
    {\iflangja{直ちに実行\\高速・ブロックしない}{runs immediately;\\fast, non-blocking}};
  \node[lab, anchor=north] at (92,-1)
    {\iflangja{スレッドとしてスケジュール\\ブロック可能}{scheduled like a thread;\\may block}};
  \draw[->] (92,20) -- node[lab, above] {\iflangja{時間}{time}} (116,20);
\end{tikzpicture}

  \caption{Interrupt handling divided into a top half, which interrupts
    thread T, and a bottom half, which the scheduler runs later.}
  \label{fig:l05-top-bottom}
\end{figure}

Handling interrupts as threads is not free.  In SunOS~5.0 it adds about 40 SPARC instructions to
every interrupt.  It also saves work elsewhere: since a handler that needs
a mutex can now block, the mutex operations no longer have to change the
interrupt mask to keep handlers from running while a mutex is held, which
saves about 12 instructions for each mutex operation.  With $c$
instructions added per interrupt and $s$ saved per mutex operation, the
scheme pays off when
\begin{equation}
  c\, N_{\mathrm{int}} \;<\; s\, N_{\mathrm{mutex}}
  \quad\Longleftrightarrow\quad
  \frac{N_{\mathrm{int}}}{N_{\mathrm{mutex}}} \;<\; \frac{s}{c},
  \label{eq:l05-break-even}
\end{equation}
where $N_{\mathrm{int}}$ and $N_{\mathrm{mutex}}$ are the numbers of
interrupts and of mutex operations.  Mutexes are locked and unlocked far
more often than interrupts occur, so the balance is
positive.\margin{Linux makes the same trade: booting with
\texttt{threadirqs} runs handlers in threads, and under
\texttt{PREEMPT\_RT}, mainline since Linux~6.12, that is the default for
every handler not marked \texttt{IRQF\_NO\_THREAD}.}

\begin{example}
  A system handles 2\,000 interrupts and performs 150\,000 mutex operations
  per second.  With $c = 40$ and $s = 12$, interrupt threads cost
  $40 \times 2\,000 = 80\,000$ instructions per second and save
  $12 \times 150\,000 = 1\,800\,000$, a net saving of $1\,720\,000$
  instructions per second.  By \cref{eq:l05-break-even} the scheme would
  lose only with more than $12/40 = 0.3$ interrupts per mutex operation,
  that is, more than 45\,000 interrupts per second at this mutex rate.
\end{example}

\begin{keypoint}
  Optimize for the common case.  Interrupts as threads make the rare event,
  an interrupt, slightly more expensive, in exchange for making the
  frequent event, a mutex operation, cheaper.
\end{keypoint}

\section{Threads and signal handling}
\label{sec:l05-thread-signals}

In a many-to-many implementation, signal masks exist at two
levels.\source{P2L4, threads and signal handling, cases~1--4; Stein and
Shah \cite{stein1992}.}  The user-level library keeps a mask for each
user-level thread in its thread data structure, and the kernel keeps a
mask for each kernel-level thread, in its LWP.  A user-level thread changes
its own mask through the library, without a system call; the kernel sees
only the masks of the kernel-level threads.  When a signal arrives, the
kernel interrupts a kernel-level thread whose mask enables it, whatever the
mask of the user-level thread currently running on it.\margin{With NPTL the
two levels coincide and the mismatch disappears: \texttt{pthread\_sigmask}
is a system call (\texttt{rt\_sigprocmask}) that sets the kernel's mask of
that thread directly.}

The SunOS library resolves the mismatch by installing its own routine as
the handler of the signals in the signal handler table of the process.
This routine runs first, examines the user-level masks, and invokes the
application's handler in a suitable user-level thread.  Four cases arise
(\cref{fig:l05-signal-cases}).
\begin{description}
  \item[Case 1] The signal is enabled both in the kernel-level thread and
    in the user-level thread running on it.  The kernel interrupts the
    kernel-level thread; the library routine finds the signal enabled in
    the current user-level thread and runs the handler there.
  \item[Case 2] The running user-level thread has the signal disabled, but
    another user-level thread, runnable but not running, has it enabled.
    The kernel interrupts the kernel-level thread, whose mask enables the
    signal.  The library routine sees that the current thread has the
    signal masked, schedules the runnable thread on this kernel-level
    thread, and the handler runs in it.
  \item[Case 3] The user-level thread with the signal enabled is running
    on a different kernel-level thread.  The library routine on the
    interrupted kernel-level thread knows where that thread runs and makes a
    system call that sends a \emph{directed} signal to the other
    kernel-level thread.  The kernel interrupts it, and the library routine
    there finds the signal enabled and runs the handler.
  \item[Case 4] No user-level thread has the signal enabled, but the masks
    of the kernel-level threads still enable it.  The library routine on
    the interrupted kernel-level thread makes a system call that disables
    the signal in the mask of that kernel-level thread, and re-raises the
    signal for the process.  The kernel delivers it to another
    kernel-level thread that still has it enabled, where the same steps
    follow, until the signal is disabled in all kernel-level threads and
    remains pending.  When a user-level thread later enables the signal,
    the library makes a system call to enable it in a kernel-level thread
    again, and the pending signal is delivered.\tip{One rule generates the
    four cases: run the handler wherever a user-level mask enables the
    signal; if no thread does, clear the kernel-level mask and re-raise.}
\end{description}

\begin{figure}[htbp]
  \centering
  % Signal handling with user-level (U) and kernel-level (K) signal masks:
% the four cases handled by the SunOS user-level thread library.
% Usage: % Signal handling with user-level (U) and kernel-level (K) signal masks:
% the four cases handled by the SunOS user-level thread library.
% Usage: % Signal handling with user-level (U) and kernel-level (K) signal masks:
% the four cases handled by the SunOS user-level thread library.
% Usage: \input{figures/l05-signal-cases}   (width ~116 mm)
\begin{tikzpicture}[x=1mm, y=1mm, font=\scriptsize\sffamily,
  >={Stealth[length=1.5mm]},
  ut/.style={draw, circle, fill=white, minimum size=5mm, inner sep=0pt},
  kt/.style={draw, rectangle, fill=ln-box, minimum size=5mm, inner sep=0pt},
  on/.style={font=\scriptsize\sffamily\bfseries, text=ln-tip, anchor=west},
  off/.style={font=\scriptsize\sffamily\bfseries, text=ln-caution, anchor=west},
  lab/.style={font=\scriptsize\sffamily, align=center},
  act/.style={->, thick, ln-accent},
  sig/.style={->, thick, ln-caution},
  ttl/.style={font=\scriptsize\sffamily, anchor=north west}]
  % \lsframe{title}: panel frame, title, user/kernel boundary, signal into K1
  \def\lsframe#1{%
    \draw[ln-rule, rounded corners=2pt] (0,-6) rectangle (56,31);
    \node[ttl] at (1,30.5) {#1};
    \draw[dashed, ln-rule] (1,12) -- (55,12);
    \draw[sig] (12,-5) -- (12,2.5);
    \node[lab, text=ln-caution, anchor=east] at (11,-2.5) {\iflangja{シグナル}{signal}};}
  % ---------------------------------------------------------------- case 1
  \begin{scope}[shift={(0,40)}]
    \lsframe{\textbf{\iflangja{ケース 1}{Case 1}}\enspace
      \iflangja{実行中の U1 が有効}{running U1 enabled}}
    \node[ut] (u1) at (12,19) {U1}; \node[on] at (14.8,19) {1};
    \node[kt] (k1) at (12,5) {K1};  \node[on] at (14.8,5) {1};
    \draw (u1) -- (k1);
    \node[lab, anchor=west, align=left] at (22,19)
      {\iflangja{ライブラリ：U1 は有効\\→ U1 でハンドラを実行}%
                {library: U1 enabled\\$\to$ handler runs in U1}};
  \end{scope}
  % ---------------------------------------------------------------- case 2
  \begin{scope}[shift={(60,40)}]
    \lsframe{\textbf{\iflangja{ケース 2}{Case 2}}\enspace
      \iflangja{実行可能な U2 が有効}{runnable U2 enabled}}
    \node[ut] (u1) at (12,19) {U1}; \node[off] at (14.8,19) {0};
    \node[kt] (k1) at (12,5) {K1};  \node[on] at (14.8,5) {1};
    \draw (u1) -- (k1);
    \node[ut, dashed] (u2) at (40,21) {U2}; \node[on] at (42.8,21) {1};
    \node[lab, text=ln-muted, anchor=north] at (40,18) {\iflangja{実行待ち}{ready}};
    \draw[act, dashed] (35,15.5) -- node[lab, below right, pos=0.5, text=ln-accent]
      {\iflangja{K1 に割り当て}{scheduled on K1}} (18.5,6.5);
  \end{scope}
  % ---------------------------------------------------------------- case 3
  \begin{scope}[shift={(0,0)}]
    \lsframe{\textbf{\iflangja{ケース 3}{Case 3}}\enspace
      \iflangja{別の K 上の U2 が有効}{U2 on another K enabled}}
    \node[ut] (u1) at (12,19) {U1}; \node[off] at (14.8,19) {0};
    \node[kt] (k1) at (12,5) {K1};  \node[on] at (14.8,5) {1};
    \draw (u1) -- (k1);
    \node[ut] (u2) at (42,19) {U2}; \node[on] at (44.8,19) {1};
    \node[kt] (k2) at (42,5) {K2};  \node[on] at (44.8,5) {1};
    \draw (u2) -- (k2);
    \draw[act] (19,5) -- node[lab, below, text=ln-accent]
      {\iflangja{指定した\\シグナル}{directed\\signal}} (39.2,5);
  \end{scope}
  % ---------------------------------------------------------------- case 4
  \begin{scope}[shift={(60,0)}]
    \lsframe{\textbf{\iflangja{ケース 4}{Case 4}}\enspace
      \iflangja{有効な U がない}{no U enabled}}
    \node[ut] (u1) at (12,19) {U1}; \node[off] at (14.8,19) {0};
    \node[kt] (k1) at (12,5) {K1};
    \node[font=\scriptsize\sffamily\bfseries, anchor=west] at (14.8,5)
      {\textcolor{ln-tip}{1}$\to$\textcolor{ln-caution}{0}};
    \draw (u1) -- (k1);
    \node[ut] (u2) at (42,19) {U2}; \node[off] at (44.8,19) {0};
    \node[kt] (k2) at (42,5) {K2};
    \node[font=\scriptsize\sffamily\bfseries, anchor=west] at (44.8,5)
      {\textcolor{ln-tip}{1}$\to$\textcolor{ln-caution}{0}};
    \draw (u2) -- (k2);
    \draw[act] (24,5) -- node[lab, below, text=ln-accent]
      {\iflangja{再送出}{re-raised}} (39.2,5);
  \end{scope}
\end{tikzpicture}
   (width ~116 mm)
\begin{tikzpicture}[x=1mm, y=1mm, font=\scriptsize\sffamily,
  >={Stealth[length=1.5mm]},
  ut/.style={draw, circle, fill=white, minimum size=5mm, inner sep=0pt},
  kt/.style={draw, rectangle, fill=ln-box, minimum size=5mm, inner sep=0pt},
  on/.style={font=\scriptsize\sffamily\bfseries, text=ln-tip, anchor=west},
  off/.style={font=\scriptsize\sffamily\bfseries, text=ln-caution, anchor=west},
  lab/.style={font=\scriptsize\sffamily, align=center},
  act/.style={->, thick, ln-accent},
  sig/.style={->, thick, ln-caution},
  ttl/.style={font=\scriptsize\sffamily, anchor=north west}]
  % \lsframe{title}: panel frame, title, user/kernel boundary, signal into K1
  \def\lsframe#1{%
    \draw[ln-rule, rounded corners=2pt] (0,-6) rectangle (56,31);
    \node[ttl] at (1,30.5) {#1};
    \draw[dashed, ln-rule] (1,12) -- (55,12);
    \draw[sig] (12,-5) -- (12,2.5);
    \node[lab, text=ln-caution, anchor=east] at (11,-2.5) {\iflangja{シグナル}{signal}};}
  % ---------------------------------------------------------------- case 1
  \begin{scope}[shift={(0,40)}]
    \lsframe{\textbf{\iflangja{ケース 1}{Case 1}}\enspace
      \iflangja{実行中の U1 が有効}{running U1 enabled}}
    \node[ut] (u1) at (12,19) {U1}; \node[on] at (14.8,19) {1};
    \node[kt] (k1) at (12,5) {K1};  \node[on] at (14.8,5) {1};
    \draw (u1) -- (k1);
    \node[lab, anchor=west, align=left] at (22,19)
      {\iflangja{ライブラリ：U1 は有効\\→ U1 でハンドラを実行}%
                {library: U1 enabled\\$\to$ handler runs in U1}};
  \end{scope}
  % ---------------------------------------------------------------- case 2
  \begin{scope}[shift={(60,40)}]
    \lsframe{\textbf{\iflangja{ケース 2}{Case 2}}\enspace
      \iflangja{実行可能な U2 が有効}{runnable U2 enabled}}
    \node[ut] (u1) at (12,19) {U1}; \node[off] at (14.8,19) {0};
    \node[kt] (k1) at (12,5) {K1};  \node[on] at (14.8,5) {1};
    \draw (u1) -- (k1);
    \node[ut, dashed] (u2) at (40,21) {U2}; \node[on] at (42.8,21) {1};
    \node[lab, text=ln-muted, anchor=north] at (40,18) {\iflangja{実行待ち}{ready}};
    \draw[act, dashed] (35,15.5) -- node[lab, below right, pos=0.5, text=ln-accent]
      {\iflangja{K1 に割り当て}{scheduled on K1}} (18.5,6.5);
  \end{scope}
  % ---------------------------------------------------------------- case 3
  \begin{scope}[shift={(0,0)}]
    \lsframe{\textbf{\iflangja{ケース 3}{Case 3}}\enspace
      \iflangja{別の K 上の U2 が有効}{U2 on another K enabled}}
    \node[ut] (u1) at (12,19) {U1}; \node[off] at (14.8,19) {0};
    \node[kt] (k1) at (12,5) {K1};  \node[on] at (14.8,5) {1};
    \draw (u1) -- (k1);
    \node[ut] (u2) at (42,19) {U2}; \node[on] at (44.8,19) {1};
    \node[kt] (k2) at (42,5) {K2};  \node[on] at (44.8,5) {1};
    \draw (u2) -- (k2);
    \draw[act] (19,5) -- node[lab, below, text=ln-accent]
      {\iflangja{指定した\\シグナル}{directed\\signal}} (39.2,5);
  \end{scope}
  % ---------------------------------------------------------------- case 4
  \begin{scope}[shift={(60,0)}]
    \lsframe{\textbf{\iflangja{ケース 4}{Case 4}}\enspace
      \iflangja{有効な U がない}{no U enabled}}
    \node[ut] (u1) at (12,19) {U1}; \node[off] at (14.8,19) {0};
    \node[kt] (k1) at (12,5) {K1};
    \node[font=\scriptsize\sffamily\bfseries, anchor=west] at (14.8,5)
      {\textcolor{ln-tip}{1}$\to$\textcolor{ln-caution}{0}};
    \draw (u1) -- (k1);
    \node[ut] (u2) at (42,19) {U2}; \node[off] at (44.8,19) {0};
    \node[kt] (k2) at (42,5) {K2};
    \node[font=\scriptsize\sffamily\bfseries, anchor=west] at (44.8,5)
      {\textcolor{ln-tip}{1}$\to$\textcolor{ln-caution}{0}};
    \draw (u2) -- (k2);
    \draw[act] (24,5) -- node[lab, below, text=ln-accent]
      {\iflangja{再送出}{re-raised}} (39.2,5);
  \end{scope}
\end{tikzpicture}
   (width ~116 mm)
\begin{tikzpicture}[x=1mm, y=1mm, font=\scriptsize\sffamily,
  >={Stealth[length=1.5mm]},
  ut/.style={draw, circle, fill=white, minimum size=5mm, inner sep=0pt},
  kt/.style={draw, rectangle, fill=ln-box, minimum size=5mm, inner sep=0pt},
  on/.style={font=\scriptsize\sffamily\bfseries, text=ln-tip, anchor=west},
  off/.style={font=\scriptsize\sffamily\bfseries, text=ln-caution, anchor=west},
  lab/.style={font=\scriptsize\sffamily, align=center},
  act/.style={->, thick, ln-accent},
  sig/.style={->, thick, ln-caution},
  ttl/.style={font=\scriptsize\sffamily, anchor=north west}]
  % \lsframe{title}: panel frame, title, user/kernel boundary, signal into K1
  \def\lsframe#1{%
    \draw[ln-rule, rounded corners=2pt] (0,-6) rectangle (56,31);
    \node[ttl] at (1,30.5) {#1};
    \draw[dashed, ln-rule] (1,12) -- (55,12);
    \draw[sig] (12,-5) -- (12,2.5);
    \node[lab, text=ln-caution, anchor=east] at (11,-2.5) {\iflangja{シグナル}{signal}};}
  % ---------------------------------------------------------------- case 1
  \begin{scope}[shift={(0,40)}]
    \lsframe{\textbf{\iflangja{ケース 1}{Case 1}}\enspace
      \iflangja{実行中の U1 が有効}{running U1 enabled}}
    \node[ut] (u1) at (12,19) {U1}; \node[on] at (14.8,19) {1};
    \node[kt] (k1) at (12,5) {K1};  \node[on] at (14.8,5) {1};
    \draw (u1) -- (k1);
    \node[lab, anchor=west, align=left] at (22,19)
      {\iflangja{ライブラリ：U1 は有効\\→ U1 でハンドラを実行}%
                {library: U1 enabled\\$\to$ handler runs in U1}};
  \end{scope}
  % ---------------------------------------------------------------- case 2
  \begin{scope}[shift={(60,40)}]
    \lsframe{\textbf{\iflangja{ケース 2}{Case 2}}\enspace
      \iflangja{実行可能な U2 が有効}{runnable U2 enabled}}
    \node[ut] (u1) at (12,19) {U1}; \node[off] at (14.8,19) {0};
    \node[kt] (k1) at (12,5) {K1};  \node[on] at (14.8,5) {1};
    \draw (u1) -- (k1);
    \node[ut, dashed] (u2) at (40,21) {U2}; \node[on] at (42.8,21) {1};
    \node[lab, text=ln-muted, anchor=north] at (40,18) {\iflangja{実行待ち}{ready}};
    \draw[act, dashed] (35,15.5) -- node[lab, below right, pos=0.5, text=ln-accent]
      {\iflangja{K1 に割り当て}{scheduled on K1}} (18.5,6.5);
  \end{scope}
  % ---------------------------------------------------------------- case 3
  \begin{scope}[shift={(0,0)}]
    \lsframe{\textbf{\iflangja{ケース 3}{Case 3}}\enspace
      \iflangja{別の K 上の U2 が有効}{U2 on another K enabled}}
    \node[ut] (u1) at (12,19) {U1}; \node[off] at (14.8,19) {0};
    \node[kt] (k1) at (12,5) {K1};  \node[on] at (14.8,5) {1};
    \draw (u1) -- (k1);
    \node[ut] (u2) at (42,19) {U2}; \node[on] at (44.8,19) {1};
    \node[kt] (k2) at (42,5) {K2};  \node[on] at (44.8,5) {1};
    \draw (u2) -- (k2);
    \draw[act] (19,5) -- node[lab, below, text=ln-accent]
      {\iflangja{指定した\\シグナル}{directed\\signal}} (39.2,5);
  \end{scope}
  % ---------------------------------------------------------------- case 4
  \begin{scope}[shift={(60,0)}]
    \lsframe{\textbf{\iflangja{ケース 4}{Case 4}}\enspace
      \iflangja{有効な U がない}{no U enabled}}
    \node[ut] (u1) at (12,19) {U1}; \node[off] at (14.8,19) {0};
    \node[kt] (k1) at (12,5) {K1};
    \node[font=\scriptsize\sffamily\bfseries, anchor=west] at (14.8,5)
      {\textcolor{ln-tip}{1}$\to$\textcolor{ln-caution}{0}};
    \draw (u1) -- (k1);
    \node[ut] (u2) at (42,19) {U2}; \node[off] at (44.8,19) {0};
    \node[kt] (k2) at (42,5) {K2};
    \node[font=\scriptsize\sffamily\bfseries, anchor=west] at (44.8,5)
      {\textcolor{ln-tip}{1}$\to$\textcolor{ln-caution}{0}};
    \draw (u2) -- (k2);
    \draw[act] (24,5) -- node[lab, below, text=ln-accent]
      {\iflangja{再送出}{re-raised}} (39.2,5);
  \end{scope}
\end{tikzpicture}

  \caption{The four cases of signal handling with user-level threads (U)
    and kernel-level threads (K).  The digit beside each thread is its mask
    bit for the signal: 1 enabled, 0 disabled.}
  \label{fig:l05-signal-cases}
\end{figure}

The design follows the same principle as interrupts as threads.  Signals
occur far less often than threads change their signal masks, typically
around every critical section.  Changing a user-level mask is a cheap
library operation, whereas keeping the kernel-level mask in step would
require a system call on every change.  In return, delivering a signal
becomes more expensive---the library routine runs first and may have to
redirect the signal or make system calls---but this happens rarely.

\section{Tasks in Linux}
\label{sec:l05-linux}

The Linux kernel does not distinguish between processes and
threads.\source{P2L4, tasks in Linux; Silberschatz et al.\ ch.~4; Kerrisk
ch.~28.}  Both are built from a single kind of execution context, and a
multithreaded process is a group of such contexts that share their
resources.

\begin{definition}[Task]
  In the Linux kernel a \term{task}[タスク] is the basic execution
  context, a kernel-level thread represented by a
  \texttt{struct task\_struct}.  A single-threaded process consists of one
  task, a multithreaded process of several tasks.
\end{definition}

\needspace{16\baselineskip}
The main fields of the structure are:
\begin{lstlisting}[language=C, numbers=none]
struct task_struct {       // abridged, Linux 3.x
    // ...
    pid_t pid;             // task identifier
    pid_t tgid;            // thread group identifier
    int prio;
    volatile long state;
    struct mm_struct *mm;         // address space
    struct files_struct *files;   // open files
    struct list_head tasks;
    int on_cpu;
    cpumask_t cpus_allowed;
    // ...
};
\end{lstlisting}
Every task has its own identifier, \texttt{pid}.  The tasks of one process
form a thread group whose identifier, \texttt{tgid}, is the \texttt{pid}
of the leader of the group; this is the process identifier that
\texttt{getpid} returns.\caution{The field named \texttt{pid} is the thread
identifier: \texttt{getpid} returns \texttt{tgid}, and \texttt{gettid}
returns \texttt{pid}.  \texttt{ps -L} prints the two as \texttt{PID} and
\texttt{LWP}.}  The task structure does not contain the
process's resources: \texttt{mm} and \texttt{files} point to separately
allocated structures for the address space and the open files.  As in
\cref{sec:l05-structures}, the state is split into structures that several
tasks can share.

A task is created with the \term{clone}[clone] system call:
\begin{lstlisting}[language=C, numbers=none]
clone(function, stack_ptr, sharing_flags, args);
\end{lstlisting}
The new task executes \texttt{function(args)} on the stack
\texttt{stack\_ptr}; each bit of \texttt{sharing\_flags} selects whether a
part of the state is shared with the creating task or copied
(\cref{tab:l05-clone}).  With all these flags set, together with
\texttt{CLONE\_THREAD}, which places the new task in the thread group of
its creator, the new task is a thread of the creating process; with all of
them cleared, it is a new process, and \texttt{fork} is implemented in
this way.

\begin{table}[htbp]
  \centering
  \caption{Sharing flags of \texttt{clone}.}
  \label{tab:l05-clone}
  \small
  \begin{tabular}{@{}l>{\raggedright\arraybackslash}p{0.33\linewidth}>{\raggedright\arraybackslash}p{0.33\linewidth}@{}}
    \toprule
    Flag & When set & When cleared \\
    \midrule
    \texttt{CLONE\_FILES}   & share the file descriptor table
                            & copy the file descriptor table \\
    \texttt{CLONE\_FS}      & share file system information
                            & copy file system information \\
    \texttt{CLONE\_SIGHAND} & share the signal handler table
                            & copy the signal handler table \\
    \texttt{CLONE\_VM}      & share the address space
                            & new address space (copy) \\
    \bottomrule
  \end{tabular}
\end{table}

POSIX threads in current Linux systems are implemented by the
\term{Native POSIX Thread Library}[NPTL]\index{NPTL|see{Native POSIX
Thread Library}} (NPTL), which uses the one-to-one model: every user-level
thread is a task.  The kernel therefore sees every thread with its
information, such as its signal mask, and the coordination problems of
\cref{sec:l05-interaction,sec:l05-thread-signals} do not arise.  The
drawbacks that once argued against the one-to-one model have faded: traps
into the kernel have become much cheaper, and memory and the range of
thread identifiers are large enough for many kernel-level threads.  NPTL
replaced LinuxThreads, an earlier library that also created one task per
thread but did not conform to POSIX in areas such as signal handling; the
many-to-many library developed at the same time, Next Generation POSIX
Threads (NGPT), had the coordination problems of
\cref{sec:l05-interaction} and was abandoned in favour of
NPTL.\margin{LinuxThreads is sometimes described as many-to-many; the
many-to-many library of that era was NGPT.}

\begin{keypoint}[Lesson summary]
  Thread state is divided into hard process state shared by all threads,
  light process state per kernel-level thread, and user-level structures
  with red zones.  In SunOS~5.0 an LWP links the two levels, which
  coordinate through system calls and signals.  Handlers run in the
  interrupted thread's context, so masks prevent deadlock; interrupts as
  threads, top and bottom halves and user-level signal masks optimize the
  common case.  Linux makes every thread a task created with
  \texttt{clone}, mapped one-to-one by NPTL.
\end{keypoint}

\end{document}
